\chapter{Partial Trace, Feedback, and Geometry of Interaction for Mealy Morphisms}
\label{ch:partial-trace-mealy}

\section{Overview}
\label{sec:partial-trace-overview}

The previous chapters developed internal categories as monads in spans,
Mealy morphisms as stateful monad morphisms, relative program categories, and
polynomial transposes of relative Mealy actions.  The present chapter shows
that this structure already contains a canonical \emph{partial trace}.

The trace construction is not added as an external axiom.  It is induced by
the labelled input-discrete presentation of Mealy morphisms.  A Mealy morphism
\[
\Phi:\mathbb A\rightsquigarrow\mathbb B
\]
is equivalently a labelled input-discrete span of internal categories
\[
\begin{tikzcd}[column sep=large]
\mathbb A^{\mathrm{op}}
&
\mathbb P
  \arrow[l, "I"']
  \arrow[r, "\Omega"]
&
\mathbb B^{\mathrm{op}},
\end{tikzcd}
\]
where the left leg \(I\) is input-discrete.  Thus Mealy morphisms admit a
faithful symmetric monoidal representation as a subcategory of the span
category of internal categories.  The ambient span category has the standard
total trace associated to spans.  The Mealy category therefore inherits a
partial trace by restriction: the trace of a Mealy morphism is defined exactly
when the ambient traced span is again input-discrete.

The restriction mechanism is the one isolated by
Malherbe--Scott--Selinger: every symmetric monoidal subcategory of a totally
traced category inherits a partial trace, and every partially traced category
embeds faithfully in a totally traced category
\cite{MalherbeScottSelingerPartiallyTracedCategories}.  The ambient trace on
spans is the standard traced monoidal structure of spans, in the sense of
Joyal--Street--Verity
\cite{JoyalStreetVerityTracedMonoidalCategories}.

The main construction is as follows.  Given a Mealy morphism
\[
\Phi:\mathbb A\times\mathbb U
\rightsquigarrow
\mathbb B\times\mathbb U,
\]
write its labelled input-discrete presentation as
\[
\begin{tikzcd}[column sep=large]
(\mathbb A\times\mathbb U)^{\mathrm{op}}
&
\mathbb P
  \arrow[l, "I"']
  \arrow[r, "\Omega"]
&
(\mathbb B\times\mathbb U)^{\mathrm{op}} .
\end{tikzcd}
\]
Using the canonical product-opposite identifications
\[
(\mathbb A\times\mathbb U)^{\mathrm{op}}
\cong
\mathbb A^{\mathrm{op}}\times\mathbb U^{\mathrm{op}},
\qquad
(\mathbb B\times\mathbb U)^{\mathrm{op}}
\cong
\mathbb B^{\mathrm{op}}\times\mathbb U^{\mathrm{op}},
\]
write
\[
I=(I_A,I_U),
\qquad
\Omega=(\Omega_B,\Omega_U).
\]
Thus
\[
I_U,\Omega_U:\mathbb P\longrightarrow\mathbb U^{\mathrm{op}}.
\]
The ambient span trace closes the hidden interface by forming the equalizer
\[
\mathbb P^\circ
:=
\operatorname{Eq}(I_U,\Omega_U).
\]
Explicitly, \(\mathbb P^\circ\) is characterized by the pullback square
\[
\begin{tikzcd}[column sep=large,row sep=large]
\mathbb P^\circ
  \arrow[r, hook, "\iota"]
  \arrow[d]
&
\mathbb P
  \arrow[d, "{(I_U,\Omega_U)}"]
\\
\mathbb U^{\mathrm{op}}
  \arrow[r, "\Delta_{\mathbb U^{\mathrm{op}}}"']
&
\mathbb U^{\mathrm{op}}\times\mathbb U^{\mathrm{op}} .
\end{tikzcd}
\]
The trace is defined precisely when
\[
I_A:\mathbb P^\circ\to\mathbb A^{\mathrm{op}}
\]
is input-discrete.  When this holds, the traced Mealy morphism is represented
by
\[
\begin{tikzcd}[column sep=large]
\mathbb A^{\mathrm{op}}
&
\mathbb P^\circ
  \arrow[l, "I_A"']
  \arrow[r, "\Omega_B"]
&
\mathbb B^{\mathrm{op}} .
\end{tikzcd}
\]

In component notation, if
\[
\Phi=(P,\delta,\omega):
\mathbb A\times\mathbb U
\rightsquigarrow
\mathbb B\times\mathbb U,
\]
the trace-class condition says that the internal feedback equation
\[
h=\omega_U((a,h),x)
\]
has a unique solution for every closed state \(x\in P^\circ\) and every
visible input arrow \(a\) admissible at \(x\), meaning
\[
t_A(a)=p_A(x).
\]
Thus the trace class is the internal unique-solvability condition saying that
feedback through \(\mathbb U\) is deterministic and remains Mealy.

The trace also has an \emph{interface-privacy} interpretation.  Tracing over
\(\mathbb U\) hides the internal \(\mathbb U\)-conversation from the visible
\(\mathbb A\)-to-\(\mathbb B\) interface.  This agrees with the interpretation
of trace as a tool for hiding what happens in the middle of an interaction,
as emphasized by Pavlovi\'c in the analysis of man-in-the-middle interactions
\cite{PavlovicTracingManInTheMiddle}.  In the present setting the hidden
middle is the feedback witness \(h\), and trace-class unique solvability is
the condition under which the hidden middle can be eliminated while retaining
a deterministic Mealy morphism.

Finally, the partial trace gives a partial Geometry-of-Interaction category of
bidirectional Mealy processes.  In the total traced case, the
Joyal--Street--Verity Int-construction yields a compact closed category
\cite{JoyalStreetVerityTracedMonoidalCategories}.  In the partial case,
Malherbe--Scott--Selinger construct the corresponding partial Int-construction
using Freyd paracategories
\cite{MalherbeScottSelingerPartiallyTracedCategories}.  Here this gives a
symmetric monoidal paracategory
\[
\operatorname{Int}_{\partial}(\mathsf{Mly}(\mathcal E))
\]
whose composition is Geometry-of-Interaction execution: form the open pipeline
and then partially trace the shared counterflow interface.  The associated
multiplicative partial linear logic has bidirectional Mealy interfaces as
formulas, bidirectional Mealy processes as proofs, and cut defined exactly
when the corresponding feedback equation is trace-class.  This is the Mealy
analogue of typed Geometry of Interaction based on partial traces and trace
classes
\cite{HaghverdiTypedGoI,AbramskyHaghverdiScottGoI}.

\[
\begin{tikzcd}[column sep=huge,row sep=large]
\Phi:\mathbb A\times\mathbb U\rightsquigarrow\mathbb B\times\mathbb U
  \arrow[r, maps to]
&
\left(
(\mathbb A\times\mathbb U)^{\mathrm{op}}
\xleftarrow{I}
\mathbb P
\xrightarrow{\Omega}
(\mathbb B\times\mathbb U)^{\mathrm{op}}
\right)
  \arrow[r, maps to]
&
\left(
\mathbb A^{\mathrm{op}}
\xleftarrow{I_A}
\mathbb P^\circ
\xrightarrow{\Omega_B}
\mathbb B^{\mathrm{op}}
\right).
\end{tikzcd}
\]

\section{Standing hypotheses and span conventions}
\label{sec:partial-trace-standing-hypotheses}

Throughout this chapter, assume that \(\mathcal E\) has finite limits.  Then
the category
\[
\mathbf{Cat}_{\mathrm{int}}(\mathcal E)
\]
of internal categories and internal functors has finite limits, computed
componentwise.  In particular, products, pullbacks, and equalizers of internal
categories exist and are created by the object-of-objects and object-of-arrows
projections.

The partial-trace construction itself uses only these finite-limit hypotheses.
Whenever we use the predicate-transformer, box, diamond, fixed-point, or
polynomial-response structures of the previous chapters, the standing
hypotheses of the corresponding earlier construction are in force.

We write
\[
\mathbf{Span}(\mathbf{Cat}_{\mathrm{int}}(\mathcal E))
\]
for the ordinary span category whose morphisms are isomorphism classes of
spans of internal categories.  Equivalently, one may work with chosen
pullbacks and read all span equalities up to the canonical apex isomorphisms.
The total trace axioms for spans are understood in this isomorphism-class
sense.

An ordinary span in this category is displayed as
\[
\begin{tikzcd}[column sep=large]
\mathbb X
&
\mathbb R
  \arrow[l, "\ell"']
  \arrow[r, "r"]
&
\mathbb Y .
\end{tikzcd}
\]
Given spans
\[
\begin{tikzcd}[column sep=large]
\mathbb X
&
\mathbb R
  \arrow[l, "\ell"']
  \arrow[r, "r"]
&
\mathbb Y
\end{tikzcd}
\qquad
\text{and}
\qquad
\begin{tikzcd}[column sep=large]
\mathbb Y
&
\mathbb S
  \arrow[l, "s"']
  \arrow[r, "u"]
&
\mathbb Z,
\end{tikzcd}
\]
their composite is represented by
\[
\begin{tikzcd}[column sep=large]
\mathbb X
&
\mathbb R\times_{\mathbb Y}\mathbb S
  \arrow[l, "\ell\pi_{\mathbb R}"']
  \arrow[r, "u\pi_{\mathbb S}"]
&
\mathbb Z,
\end{tikzcd}
\]
where the apex is defined by the pullback
\[
\begin{tikzcd}[column sep=large,row sep=large]
\mathbb R\times_{\mathbb Y}\mathbb S
  \arrow[r, "\pi_{\mathbb S}"]
  \arrow[d, "\pi_{\mathbb R}"']
&
\mathbb S
  \arrow[d, "s"]
\\
\mathbb R
  \arrow[r, "r"']
&
\mathbb Y .
\end{tikzcd}
\]

All references to input-discreteness are made in the sense fixed in the
previous chapters.  Thus an internal functor
\[
J:\mathbb X\to\mathbb A^{\mathrm{op}}
\]
is \emph{input-discrete} when the square
\[
\begin{tikzcd}[column sep=large,row sep=large]
X_1
  \arrow[r,"J_1"]
  \arrow[d,"s_X"']
&
A_1
  \arrow[d,"t_A"]
\\
X_0
  \arrow[r,"J_0"']
&
A_0
\end{tikzcd}
\]
is a pullback.  This convention is adapted to right internal actions: an
arrow \((a,x)\) starts at \(x\), where \(p(x)=t_A(a)\), and executes backwards
along \(a\) to the state over \(s_A(a)\).

\section{The Mealy category inside spans of internal categories}
\label{sec:mealy-category-inside-spans-internal-categories}

We write
\[
\mathsf{Mly}(\mathcal E)
\]
for the ordinary category whose objects are internal categories in
\(\mathcal E\) and whose morphisms are Mealy morphisms, taken up to the
canonical isomorphisms of carrier spans generated by chosen pullbacks.
Composition and identities are those constructed in
Chapter~\ref{ch:spans-internal-categories-mealy}.  Tensor product is given by
product of internal categories and by componentwise tensor of Mealy morphisms.

\begin{proposition}
\label{prop:mly-symmetric-monoidal}
The category
\[
\mathsf{Mly}(\mathcal E)
\]
is symmetric monoidal under product of internal categories.
\end{proposition}

\begin{proof}
Products of internal categories are computed componentwise:
\[
(\mathbb A\times\mathbb B)_0=A_0\times B_0,
\qquad
(\mathbb A\times\mathbb B)_1=A_1\times B_1,
\]
with source, target, identities, and composition defined componentwise.

If
\[
\Phi=(P,\delta_\Phi,\omega_\Phi):
\mathbb A\rightsquigarrow\mathbb B
\]
and
\[
\Psi=(Q,\delta_\Psi,\omega_\Psi):
\mathbb C\rightsquigarrow\mathbb D,
\]
then
\[
\Phi\times\Psi:
\mathbb A\times\mathbb C
\rightsquigarrow
\mathbb B\times\mathbb D
\]
has carrier \(P\times Q\), transition map
\[
\delta_{\Phi\times\Psi}
=
\delta_\Phi\times\delta_\Psi,
\]
and output map
\[
\omega_{\Phi\times\Psi}
=
\omega_\Phi\times\omega_\Psi.
\]
The Mealy unit and multiplication laws hold componentwise.  The associativity,
unit, and symmetry constraints are inherited from products of internal
categories and from the span calculus.  Hence
\(\mathsf{Mly}(\mathcal E)\) is symmetric monoidal.
\end{proof}

Recall that a Mealy morphism
\[
\Phi:\mathbb A\rightsquigarrow\mathbb B
\]
is equivalently a labelled input-discrete internal category
\[
\begin{tikzcd}[column sep=large]
&
\mathbb P
  \arrow[dl, "I"']
  \arrow[dr, "\Omega"]
\\
\mathbb A^{\mathrm{op}}
&&
\mathbb B^{\mathrm{op}},
\end{tikzcd}
\]
with \(I\) input-discrete.  Equivalently, it is a span of internal categories
\[
\begin{tikzcd}[column sep=large]
\mathbb A^{\mathrm{op}}
&
\mathbb P
  \arrow[l, "I"']
  \arrow[r, "\Omega"]
&
\mathbb B^{\mathrm{op}}
\end{tikzcd}
\]
whose left leg is input-discrete.

\begin{definition}[The span subcategory of Mealy morphisms]
\label{def:span-embedding-mealy}
Let
\[
\mathcal S_{\mathrm{Mly}}
\subseteq
\mathbf{Span}(\mathbf{Cat}_{\mathrm{int}}(\mathcal E))
\]
be the subcategory whose objects are internal categories of the form
\(\mathbb A^{\mathrm{op}}\), and whose morphisms are spans
\[
\begin{tikzcd}[column sep=large]
\mathbb A^{\mathrm{op}}
&
\mathbb P
  \arrow[l, "I"']
  \arrow[r, "\Omega"]
&
\mathbb B^{\mathrm{op}}
\end{tikzcd}
\]
with \(I\) input-discrete, taken up to labelled isomorphism of apex internal
categories.
\end{definition}

\begin{theorem}[Mealy morphisms as input-discrete spans]
\label{thm:mealy-as-span-subcategory}
The labelled input-discrete construction gives a symmetric monoidal
equivalence of ordinary categories
\[
\mathsf{Mly}(\mathcal E)
\simeq
\mathcal S_{\mathrm{Mly}}.
\]
Composed with the inclusion
\[
\mathcal S_{\mathrm{Mly}}
\hookrightarrow
\mathbf{Span}(\mathbf{Cat}_{\mathrm{int}}(\mathcal E)),
\]
this gives a faithful symmetric monoidal span representation of Mealy
morphisms.
\end{theorem}

\begin{proof}
The equivalence on hom-sets is the labelled input-discrete characterization of
Mealy morphisms established in
Chapter~\ref{ch:relative-mealy-program-categories}.  The quotient by labelled
isomorphism of apex internal categories matches the quotient by canonical
carrier-span isomorphisms in \(\mathsf{Mly}(\mathcal E)\).

Identities correspond to identity input-discrete spans.  Composition
corresponds to pullback of labelled input-discrete spans.

Given spans
\[
\begin{tikzcd}[column sep=large]
\mathbb A^{\mathrm{op}}
&
\mathbb P
  \arrow[l, "I"']
  \arrow[r, "\Omega"]
&
\mathbb B^{\mathrm{op}}
\end{tikzcd}
\]
and
\[
\begin{tikzcd}[column sep=large]
\mathbb B^{\mathrm{op}}
&
\mathbb Q
  \arrow[l, "J"']
  \arrow[r, "\Theta"]
&
\mathbb C^{\mathrm{op}},
\end{tikzcd}
\]
their composite has apex
\[
\mathbb P\times_{\mathbb B^{\mathrm{op}}}\mathbb Q,
\]
defined by the pullback
\[
\begin{tikzcd}[column sep=large,row sep=large]
\mathbb P\times_{\mathbb B^{\mathrm{op}}}\mathbb Q
  \arrow[r, "\pi_{\mathbb Q}"]
  \arrow[d, "\pi_{\mathbb P}"']
&
\mathbb Q
  \arrow[d, "J"]
\\
\mathbb P
  \arrow[r, "\Omega"']
&
\mathbb B^{\mathrm{op}} .
\end{tikzcd}
\]
The left leg of the composite is
\[
I\pi_{\mathbb P}:
\mathbb P\times_{\mathbb B^{\mathrm{op}}}\mathbb Q
\longrightarrow
\mathbb A^{\mathrm{op}}.
\]
The input-discrete square for \(I\pi_{\mathbb P}\) is obtained by pasting the
input-discrete square for \(I\) with the pullback square defining the composite
and the input-discrete square for \(J\).  Since pullbacks compose, the
composite left leg is input-discrete.

The construction is compatible with products.  The product of two
input-discrete functors is input-discrete, since the defining pullback squares
are computed componentwise.  Therefore the equivalence is symmetric monoidal.
\end{proof}

\section{The ambient total trace on spans of internal categories}
\label{sec:ambient-total-trace-spans-internal-categories}

The ambient category
\[
\mathbf{Span}(\mathbf{Cat}_{\mathrm{int}}(\mathcal E))
\]
has the standard \emph{total trace} of spans with respect to the product
monoidal structure.

Let
\[
\begin{tikzcd}[column sep=large]
\mathbb X\times\mathbb U
&
\mathbb R
  \arrow[l, "\ell"']
  \arrow[r, "r"]
&
\mathbb Y\times\mathbb U
\end{tikzcd}
\]
be a span of internal categories.  Write
\[
\ell=(\ell_X,\ell_U),
\qquad
r=(r_Y,r_U).
\]
The trace over \(\mathbb U\) is the span
\[
\begin{tikzcd}[column sep=large]
\mathbb X
&
\mathbb R^\circ
  \arrow[l, "\ell_X\iota"']
  \arrow[r, "r_Y\iota"]
&
\mathbb Y,
\end{tikzcd}
\]
where
\[
\mathbb R^\circ
:=
\operatorname{Eq}(\ell_U,r_U),
\]
with inclusion \(\iota:\mathbb R^\circ\hookrightarrow\mathbb R\).  Equivalently,
\[
\mathbb R^\circ
=
\mathbb R\times_{\mathbb U\times\mathbb U}\mathbb U,
\]
using
\[
(\ell_U,r_U):\mathbb R\to\mathbb U\times\mathbb U
\]
and the diagonal
\[
\Delta_{\mathbb U}:\mathbb U\to\mathbb U\times\mathbb U.
\]
Thus \(\mathbb R^\circ\) is characterized by
\[
\begin{tikzcd}[column sep=large,row sep=large]
\mathbb R^\circ
  \arrow[r, hook, "\iota"]
  \arrow[d]
&
\mathbb R
  \arrow[d, "{(\ell_U,r_U)}"]
\\
\mathbb U
  \arrow[r, "\Delta_{\mathbb U}"']
&
\mathbb U\times\mathbb U .
\end{tikzcd}
\]

\begin{proposition}[Total trace of spans]
\label{prop:ambient-span-total-trace}
The operation above is the standard total trace on
\[
\mathbf{Span}(\mathbf{Cat}_{\mathrm{int}}(\mathcal E))
\]
with respect to the product monoidal structure.
\end{proposition}

\begin{proof}
The span category of any category with finite limits is compact closed under
the product monoidal structure.  The dual of an object \(\mathbb U\) is
\(\mathbb U\) itself.  The coevaluation and evaluation are represented by the
diagonal span
\[
\mathbf 1
\longleftarrow
\mathbb U
\xrightarrow{\Delta_{\mathbb U}}
\mathbb U\times\mathbb U
\]
and its opposite-direction counterpart
\[
\mathbb U\times\mathbb U
\xleftarrow{\Delta_{\mathbb U}}
\mathbb U
\longrightarrow
\mathbf 1.
\]
The trace of a span
\[
\mathbb X\times\mathbb U
\xleftarrow{\ell}
\mathbb R
\xrightarrow{r}
\mathbb Y\times\mathbb U
\]
is obtained by composing with these evaluation and coevaluation spans.  The
resulting apex is the pullback that enforces the equality of the two hidden
\(\mathbb U\)-coordinates.  Concretely, the composite apex is canonically
isomorphic to
\[
\mathbb R\times_{\mathbb U\times\mathbb U}\mathbb U,
\]
where
\[
\mathbb R\to\mathbb U\times\mathbb U
\]
is \((\ell_U,r_U)\) and
\[
\mathbb U\to\mathbb U\times\mathbb U
\]
is \(\Delta_{\mathbb U}\).  This is precisely
\[
\operatorname{Eq}(\ell_U,r_U).
\]
The trace axioms are the yanking, sliding, superposing, and vanishing
identities of the compact closed trace.  In span form, each is verified by
the universal property of the corresponding pullback and diagonal.  Hence the
displayed operation is the standard total trace.
\end{proof}

\section{Trace class for Mealy morphisms}
\label{sec:trace-class-mealy-morphisms}

Let
\[
\Phi:
\mathbb A\times\mathbb U
\rightsquigarrow
\mathbb B\times\mathbb U
\]
be a Mealy morphism.  Present it as a labelled input-discrete span
\[
\begin{tikzcd}[column sep=large]
(\mathbb A\times\mathbb U)^{\mathrm{op}}
&
\mathbb P
  \arrow[l, "I"']
  \arrow[r, "\Omega"]
&
(\mathbb B\times\mathbb U)^{\mathrm{op}} .
\end{tikzcd}
\]
Under the canonical identifications
\[
(\mathbb A\times\mathbb U)^{\mathrm{op}}
\cong
\mathbb A^{\mathrm{op}}\times\mathbb U^{\mathrm{op}},
\qquad
(\mathbb B\times\mathbb U)^{\mathrm{op}}
\cong
\mathbb B^{\mathrm{op}}\times\mathbb U^{\mathrm{op}},
\]
write
\[
I=(I_A,I_U),
\qquad
\Omega=(\Omega_B,\Omega_U).
\]
Thus
\[
I_U,\Omega_U:\mathbb P\to\mathbb U^{\mathrm{op}}.
\]

\begin{definition}[Closed-loop category]
\label{def:closed-loop-category}
The \emph{closed-loop category} of \(\Phi\) over \(\mathbb U\) is
\[
\mathbb P^\circ
:=
\operatorname{Eq}(I_U,\Omega_U).
\]
Thus \(\mathbb P^\circ\) is the internal subcategory of \(\mathbb P\) whose
hidden \(\mathbb U\)-input and hidden \(\mathbb U\)-output labels agree.
\end{definition}

The defining square is
\[
\begin{tikzcd}[column sep=large,row sep=large]
\mathbb P^\circ
  \arrow[r, hook, "\iota"]
  \arrow[d]
&
\mathbb P
  \arrow[d, "{(I_U,\Omega_U)}"]
\\
\mathbb U^{\mathrm{op}}
  \arrow[r, "\Delta_{\mathbb U^{\mathrm{op}}}"']
&
\mathbb U^{\mathrm{op}}\times\mathbb U^{\mathrm{op}} .
\end{tikzcd}
\]

The ambient span trace of \(\Phi\) is the span
\[
\begin{tikzcd}[column sep=large]
\mathbb A^{\mathrm{op}}
&
\mathbb P^\circ
  \arrow[l, "I_A\iota"']
  \arrow[r, "\Omega_B\iota"]
&
\mathbb B^{\mathrm{op}} .
\end{tikzcd}
\]
We suppress \(\iota\) from the notation when no confusion can arise.

\begin{definition}[Trace-class Mealy morphism]
\label{def:trace-class-mealy}
The Mealy morphism
\[
\Phi:
\mathbb A\times\mathbb U
\rightsquigarrow
\mathbb B\times\mathbb U
\]
is \emph{trace-class over \(\mathbb U\)} if
\[
I_A:\mathbb P^\circ\to\mathbb A^{\mathrm{op}}
\]
is input-discrete.
\end{definition}

\begin{definition}[Partial trace of a trace-class Mealy morphism]
\label{def:partial-trace-mealy}
If \(\Phi\) is trace-class over \(\mathbb U\), its \emph{partial trace} is the
Mealy morphism
\[
\operatorname{Tr}^{\mathbb U}(\Phi):
\mathbb A\rightsquigarrow\mathbb B
\]
represented by the labelled input-discrete span
\[
\begin{tikzcd}[column sep=large]
\mathbb A^{\mathrm{op}}
&
\mathbb P^\circ
  \arrow[l, "I_A"']
  \arrow[r, "\Omega_B"]
&
\mathbb B^{\mathrm{op}} .
\end{tikzcd}
\]
Equivalently, this is the ambient span trace over
\(\mathbb U^{\mathrm{op}}\), transported through the opposite convention used
in the Mealy-span representation.
\end{definition}

\begin{proposition}[Trace class as closure under ambient trace]
\label{prop:trace-class-ambient-closure}
For a Mealy morphism
\[
\Phi:
\mathbb A\times\mathbb U
\rightsquigarrow
\mathbb B\times\mathbb U,
\]
the following are equivalent:
\begin{enumerate}
\item \(\Phi\) is trace-class over \(\mathbb U\);
\item the ambient span trace of the labelled input-discrete span associated to
\(\Phi\) lies again in the Mealy subcategory
\[
\mathcal S_{\mathrm{Mly}}
\subseteq
\mathbf{Span}(\mathbf{Cat}_{\mathrm{int}}(\mathcal E)).
\]
\end{enumerate}
\end{proposition}

\begin{proof}
The ambient span trace is
\[
\begin{tikzcd}[column sep=large]
\mathbb A^{\mathrm{op}}
&
\mathbb P^\circ
  \arrow[l, "I_A"']
  \arrow[r, "\Omega_B"]
&
\mathbb B^{\mathrm{op}} .
\end{tikzcd}
\]
By definition of \(\mathcal S_{\mathrm{Mly}}\), this span lies in the Mealy
subcategory exactly when its left leg \(I_A\) is input-discrete.  This is
precisely the trace-class condition.
\end{proof}

\section{The internal feedback equation}
\label{sec:internal-feedback-equation}

We now unpack the trace-class condition in the original component notation of
Mealy morphisms.

Let
\[
\Phi=(P,\delta,\omega):
\mathbb A\times\mathbb U
\rightsquigarrow
\mathbb B\times\mathbb U.
\]
The carrier span has the form
\[
\begin{tikzcd}[column sep=large]
A_0\times U_0
&
P
  \arrow[l, "p"']
  \arrow[r, "q"]
&
B_0\times U_0 .
\end{tikzcd}
\]
Write
\[
p=(p_A,p_U),
\qquad
q=(q_B,q_U).
\]
The output component decomposes as
\[
\omega=(\omega_B,\omega_U).
\]

\begin{definition}[Closed state object]
\label{def:closed-state-object}
The \emph{closed state object} of \(\Phi\) over \(\mathbb U\) is
\[
P^\circ:=\operatorname{Eq}(p_U,q_U).
\]
A generalized element of \(P^\circ\) is a state \(x\in P\) satisfying
\[
p_U(x)=q_U(x).
\]
\end{definition}

Equivalently, \(P^\circ\) is given by the pullback
\[
\begin{tikzcd}[column sep=large,row sep=large]
P^\circ
  \arrow[r, hook, "\iota_P"]
  \arrow[d]
&
P
  \arrow[d, "{(p_U,q_U)}"]
\\
U_0
  \arrow[r, "\Delta_{U_0}"']
&
U_0\times U_0 .
\end{tikzcd}
\]

Let
\[
D_\Phi
:=
A_1\times_{t_A,A_0,p_A}P^\circ.
\]
A generalized element of \(D_\Phi\) is an admissible visible input/state pair
\[
(a,x)
\]
with
\[
t_A(a)=p_A(x).
\]

Define
\[
H_\Phi
:=
D_\Phi\times_{p_U\pi_P,U_0,t_U}U_1.
\]
A generalized element of \(H_\Phi\) is a triple
\[
(a,x,h)
\]
such that
\[
t_A(a)=p_A(x),
\qquad
t_U(h)=p_U(x)=q_U(x).
\]

There are two maps
\[
H_\Phi\to U_1.
\]
The first is the projection
\[
\pi_U:H_\Phi\to U_1,
\qquad
(a,x,h)\longmapsto h.
\]
The second is induced by the hidden output component:
\[
\chi_\Phi:H_\Phi\to U_1,
\qquad
(a,x,h)\longmapsto\omega_U((a,h),x).
\]
This map is well-defined because \((a,h)\) is an admissible arrow of
\(\mathbb A\times\mathbb U\) at the state \(x\), and because
\[
t_U\omega_U((a,h),x)=q_U(x)=p_U(x)=t_U(h).
\]

\begin{definition}[Feedback witness object]
\label{def:feedback-witness-object}
The \emph{feedback witness object} of \(\Phi\) over \(\mathbb U\) is the
equalizer
\[
W_\Phi
:=
\operatorname{Eq}
\left(
\pi_U,\chi_\Phi
\right)
\hookrightarrow
H_\Phi.
\]
Thus a generalized element of \(W_\Phi\) is a triple
\[
(a,x,h)
\]
with
\[
t_A(a)=p_A(x),
\qquad
t_U(h)=p_U(x)=q_U(x),
\qquad
h=\omega_U((a,h),x).
\]
There is a canonical projection
\[
\pi_\Phi:W_\Phi\to D_\Phi,
\qquad
(a,x,h)\longmapsto(a,x).
\]
\end{definition}

The defining diagram for \(W_\Phi\) is
\[
\begin{tikzcd}[column sep=large,row sep=large]
W_\Phi
  \arrow[r, hook, "\iota_W"]
  \arrow[dr, "\pi_\Phi"']
&
H_\Phi
  \arrow[r, shift left=0.7ex, "\pi_U"]
  \arrow[r, shift right=0.7ex, "\chi_\Phi"']
  \arrow[d, "\pi_D"]
&
U_1
\\
&
D_\Phi
&
\end{tikzcd}
\]
where \(\pi_\Phi=\pi_D\iota_W\).

\begin{proposition}[Internal trace-class criterion]
\label{prop:internal-trace-class-criterion}
For
\[
\Phi:
\mathbb A\times\mathbb U
\rightsquigarrow
\mathbb B\times\mathbb U,
\]
the following are equivalent:
\begin{enumerate}
\item \(\Phi\) is trace-class over \(\mathbb U\);
\item the functor
\[
I_A:\mathbb P^\circ\to\mathbb A^{\mathrm{op}}
\]
is input-discrete;
\item the projection
\[
\pi_\Phi:W_\Phi\to D_\Phi
\]
is an isomorphism.
\end{enumerate}
\end{proposition}

\begin{proof}
The equivalence between (1) and (2) is
Definition~\ref{def:trace-class-mealy}.

It remains to identify input-discreteness of
\[
I_A:\mathbb P^\circ\to\mathbb A^{\mathrm{op}}
\]
with the assertion that
\[
\pi_\Phi:W_\Phi\to D_\Phi
\]
is an isomorphism.

Since
\[
I:\mathbb P\to(\mathbb A\times\mathbb U)^{\mathrm{op}}
\]
is input-discrete, the square
\[
\begin{tikzcd}[column sep=large,row sep=large]
P_1
  \arrow[r, "I_1"]
  \arrow[d, "s_P"']
&
A_1\times U_1
  \arrow[d, "t_A\times t_U"]
\\
P_0
  \arrow[r, "p_A\times p_U"']
&
A_0\times U_0
\end{tikzcd}
\]
is a pullback.  Hence an arrow of \(\mathbb P\) with source \(x\) over the
input arrow \((a,h)\) is uniquely determined whenever
\[
t_A(a)=p_A(x),
\qquad
t_U(h)=p_U(x).
\]

The closed-loop category
\[
\mathbb P^\circ=\operatorname{Eq}(I_U,\Omega_U)
\]
has object object \(P^\circ=\operatorname{Eq}(p_U,q_U)\).  Its arrows are
exactly those arrows of \(\mathbb P\) whose hidden input and hidden output
labels agree.  Under the input-discrete identification of \(P_1\) above, such
an arrow is equivalently a triple
\[
(a,x,h)
\]
with
\[
t_A(a)=p_A(x),
\qquad
t_U(h)=p_U(x)=q_U(x),
\qquad
h=\omega_U((a,h),x).
\]
This is precisely a generalized element of \(W_\Phi\).

The defining input-discreteness square for
\[
I_A:\mathbb P^\circ\to\mathbb A^{\mathrm{op}}
\]
has comparison map
\[
P^\circ_1
\longrightarrow
A_1\times_{t_A,A_0,p_A}P^\circ
=
D_\Phi.
\]
Under the identification \(P^\circ_1\cong W_\Phi\), this comparison map is
exactly
\[
\pi_\Phi:W_\Phi\to D_\Phi.
\]
Therefore \(I_A\) is input-discrete if and only if \(\pi_\Phi\) is an
isomorphism.
\end{proof}

The ambient span trace keeps all hidden witnesses satisfying the feedback
equation.  The Mealy trace is defined exactly when this hidden-witness
relation is functional over visible admissible inputs and closed states.

When \(\Phi\) is trace-class, denote the unique feedback solution by
\[
h_\Phi(a,x).
\]

\begin{proposition}[Component formula for the partial trace]
\label{prop:component-formula-partial-trace}
Let
\[
\Phi=(P,\delta,\omega):
\mathbb A\times\mathbb U
\rightsquigarrow
\mathbb B\times\mathbb U
\]
be trace-class over \(\mathbb U\).  Then
\[
\operatorname{Tr}^{\mathbb U}(\Phi):
\mathbb A\rightsquigarrow\mathbb B
\]
has carrier span
\[
\begin{tikzcd}[column sep=large]
A_0
&
P^\circ
  \arrow[l, "p_A"']
  \arrow[r, "q_B"]
&
B_0,
\end{tikzcd}
\]
and its transition and output maps are
\[
\delta^{\operatorname{tr}}(a,x)
=
\delta((a,h_\Phi(a,x)),x),
\]
and
\[
\omega^{\operatorname{tr}}(a,x)
=
\omega_B((a,h_\Phi(a,x)),x).
\]
\end{proposition}

\begin{proof}
By Definition~\ref{def:partial-trace-mealy}, the traced Mealy morphism is
represented by the input-discrete span
\[
\begin{tikzcd}[column sep=large]
\mathbb A^{\mathrm{op}}
&
\mathbb P^\circ
  \arrow[l, "I_A"']
  \arrow[r, "\Omega_B"]
&
\mathbb B^{\mathrm{op}} .
\end{tikzcd}
\]
The object part of this span is
\[
\begin{tikzcd}[column sep=large]
A_0
&
P^\circ
  \arrow[l, "p_A"']
  \arrow[r, "q_B"]
&
B_0 .
\end{tikzcd}
\]
The arrow part is determined by the unique arrow of \(\mathbb P^\circ\) over a
visible input arrow \(a\) at a closed state \(x\).  By
Proposition~\ref{prop:internal-trace-class-criterion}, this arrow is
equivalently the unique feedback witness \(h_\Phi(a,x)\).  Substituting this
witness into the original Mealy transition and output maps gives the displayed
formulas.

It remains to check that the displayed transition lands in \(P^\circ\).  Let
\[
h=h_\Phi(a,x).
\]
The Mealy typing equations give
\[
p_U\delta((a,h),x)=s_U(h),
\]
and
\[
q_U\delta((a,h),x)=s_U\omega_U((a,h),x).
\]
Since
\[
h=\omega_U((a,h),x),
\]
these two objects are equal.  Hence
\[
\delta((a,h),x)\in P^\circ.
\]
\end{proof}

\section{The partially traced Mealy category}
\label{sec:partially-traced-mealy-category}

\begin{theorem}[Partial trace on Mealy morphisms]
\label{thm:partial-trace-on-mealy}
The symmetric monoidal category
\[
\mathsf{Mly}(\mathcal E)
\]
carries a canonical partial trace
\[
\operatorname{Tr}^{\mathbb U}_{\mathbb A,\mathbb B}:
\mathsf{Mly}(\mathcal E)(\mathbb A\times\mathbb U,\mathbb B\times\mathbb U)
\rightharpoonup
\mathsf{Mly}(\mathcal E)(\mathbb A,\mathbb B).
\]
The trace is defined exactly on trace-class Mealy morphisms in the sense of
Definition~\ref{def:trace-class-mealy}.
\end{theorem}

\begin{proof}
By Theorem~\ref{thm:mealy-as-span-subcategory}, the category
\(\mathsf{Mly}(\mathcal E)\) is symmetrically monoidally equivalent to the
subcategory
\[
\mathcal S_{\mathrm{Mly}}
\subseteq
\mathbf{Span}(\mathbf{Cat}_{\mathrm{int}}(\mathcal E))
\]
whose morphisms are spans with input-discrete left leg.

By Proposition~\ref{prop:ambient-span-total-trace}, the ambient span category
is totally traced.  By the restriction theorem of
Malherbe--Scott--Selinger, every symmetric monoidal subcategory of a totally
traced category inherits a partial trace by restriction
\cite{MalherbeScottSelingerPartiallyTracedCategories}.  Thus
\(\mathcal S_{\mathrm{Mly}}\), and hence \(\mathsf{Mly}(\mathcal E)\), is
partially traced.

The inherited trace is defined exactly when the ambient traced span lies again
in the subcategory.  By Proposition~\ref{prop:trace-class-ambient-closure},
this is exactly the trace-class condition.
\end{proof}

The restriction diagram is
\[
\begin{tikzcd}[column sep=huge,row sep=large]
\mathcal S_{\mathrm{Mly}}(\mathbb A^{\mathrm{op}}\times\mathbb U^{\mathrm{op}},
\mathbb B^{\mathrm{op}}\times\mathbb U^{\mathrm{op}})
  \arrow[r, hook]
  \arrow[d, dashed, "{\operatorname{Tr}^{\mathbb U}}"']
&
\mathbf{Span}(\mathbf{Cat}_{\mathrm{int}}(\mathcal E))
(\mathbb A^{\mathrm{op}}\times\mathbb U^{\mathrm{op}},
\mathbb B^{\mathrm{op}}\times\mathbb U^{\mathrm{op}})
  \arrow[d, "{\operatorname{Tr}^{\mathbb U^{\mathrm{op}}}}"]
\\
\mathcal S_{\mathrm{Mly}}(\mathbb A^{\mathrm{op}},\mathbb B^{\mathrm{op}})
  \arrow[r, hook]
&
\mathbf{Span}(\mathbf{Cat}_{\mathrm{int}}(\mathcal E))
(\mathbb A^{\mathrm{op}},\mathbb B^{\mathrm{op}}).
\end{tikzcd}
\]
The dashed arrow is defined on precisely those morphisms whose ambient trace
belongs to the lower-left hom-set.

\begin{corollary}[Embedding into a totally traced category]
\label{cor:mealy-partial-trace-total-embedding}
The partially traced category
\[
\mathsf{Mly}(\mathcal E)
\]
embeds faithfully into a totally traced category.  Concretely, the embedding is
the labelled input-discrete span representation
\[
\mathsf{Mly}(\mathcal E)
\hookrightarrow
\mathbf{Span}(\mathbf{Cat}_{\mathrm{int}}(\mathcal E)).
\]
It is trace-preserving and trace-reflecting on trace-class morphisms.
\end{corollary}

\begin{proof}
The concrete embedding is the embedding of
Theorem~\ref{thm:mealy-as-span-subcategory}.  The ambient category is totally
traced by Proposition~\ref{prop:ambient-span-total-trace}.  The representation
theorem of Malherbe--Scott--Selinger states that every partially traced
category embeds faithfully in a totally traced category, and their restriction
theorem identifies partial traces on symmetric monoidal subcategories with
restricted total traces
\cite{MalherbeScottSelingerPartiallyTracedCategories}.  In the present case,
the concrete span embedding realizes this representation.
\end{proof}

\section{Trace as privacy of the hidden interface}
\label{sec:trace-as-privacy-hidden-interface}

The partial trace has an interface-privacy reading.  Let
\[
\Phi:
\mathbb A\times\mathbb U
\rightsquigarrow
\mathbb B\times\mathbb U
\]
be trace-class over \(\mathbb U\).  The untraced system exposes both the
visible interface
\[
\mathbb A\rightsquigarrow\mathbb B
\]
and the hidden interface
\[
\mathbb U\rightsquigarrow\mathbb U.
\]
The traced system
\[
\operatorname{Tr}^{\mathbb U}(\Phi):
\mathbb A\rightsquigarrow\mathbb B
\]
exposes only the visible interface.  The hidden feedback witness
\[
h_\Phi(a,x)
\]
is used internally to compute the transition and output, but it is not a
label in the traced Mealy span.

This is the categorical form of hiding an internal interaction.  Pavlovi\'c
uses traced monoidal categories to model man-in-the-middle interactions and
emphasizes that trace hides what happens in the middle of a communication
scenario
\cite{PavlovicTracingManInTheMiddle}.  In the present setting the hidden
middle is the \(\mathbb U\)-feedback channel.  The trace-class condition says
that this hidden middle is internally determined by the visible input and the
closed state.

\begin{definition}[Visible observation]
\label{def:visible-observation}
Let
\[
\Phi:
\mathbb A\times\mathbb U
\rightsquigarrow
\mathbb B\times\mathbb U
\]
be trace-class.  A \emph{visible state observation} of the closed-loop
behaviour is a predicate
\[
\theta\hookrightarrow P^\circ_0
\]
for which there exists a predicate
\[
\bar\theta\hookrightarrow A_0\times B_0
\]
such that
\[
\theta=(I_{A,0},\Omega_{B,0})^*\bar\theta.
\]
A \emph{visible transition observation} is a predicate
\[
\psi\hookrightarrow P^\circ_1
\]
for which there exists a predicate
\[
\bar\psi\hookrightarrow A_1\times B_1
\]
such that
\[
\psi=(I_{A,1},\Omega_{B,1})^*\bar\psi.
\]
A \emph{visible functorial observation} is a functor
\[
\mathcal O:\mathbb P^\circ\to\mathbb V
\]
that factors through the visible labelling
\[
(I_A,\Omega_B):
\mathbb P^\circ\to
\mathbb A^{\mathrm{op}}\times\mathbb B^{\mathrm{op}}.
\]
\end{definition}

For functorial observations, visibility means that there exists a functor
\[
\overline{\mathcal O}:
\mathbb A^{\mathrm{op}}\times\mathbb B^{\mathrm{op}}\to\mathbb V
\]
such that
\[
\mathcal O=
\overline{\mathcal O}(I_A,\Omega_B).
\]
This is the commutative diagram
\[
\begin{tikzcd}[column sep=large,row sep=large]
\mathbb P^\circ
  \arrow[rr, "\mathcal O"]
  \arrow[dr, "{(I_A,\Omega_B)}"']
&&
\mathbb V
\\
&
\mathbb A^{\mathrm{op}}\times\mathbb B^{\mathrm{op}}
  \arrow[ur, "\overline{\mathcal O}"']
&
\end{tikzcd}
\]

\begin{proposition}[Visible observations factor through the traced span]
\label{prop:privacy-by-feedback-hiding}
Let \(\Phi\) be trace-class over \(\mathbb U\).  Then every visible state,
transition, or functorial observation of the closed-loop behaviour of \(\Phi\)
factors through the traced Mealy morphism
\[
\operatorname{Tr}^{\mathbb U}(\Phi).
\]
\end{proposition}

\begin{proof}
By Definition~\ref{def:partial-trace-mealy}, the traced Mealy morphism is
represented by
\[
\begin{tikzcd}[column sep=large]
\mathbb A^{\mathrm{op}}
&
\mathbb P^\circ
  \arrow[l, "I_A"']
  \arrow[r, "\Omega_B"]
&
\mathbb B^{\mathrm{op}} .
\end{tikzcd}
\]
This is exactly the closed-loop labelled category with the hidden
\(\mathbb U\)-component omitted.  State observations factor through
\[
(I_{A,0},\Omega_{B,0}):P^\circ_0\to A_0\times B_0,
\]
transition observations factor through
\[
(I_{A,1},\Omega_{B,1}):P^\circ_1\to A_1\times B_1,
\]
and functorial observations factor through
\[
(I_A,\Omega_B):
\mathbb P^\circ\to
\mathbb A^{\mathrm{op}}\times\mathbb B^{\mathrm{op}}.
\]
These are precisely the visible components of the traced span.
\end{proof}

\section{Dynamic logic of feedback closure}
\label{sec:dynamic-logic-feedback-closure}

We now express the partial trace in the dynamic logic developed in
Chapter~\ref{ch:relative-mealy-program-categories}.

Let
\[
\Phi:
\mathbb A\times\mathbb U
\rightsquigarrow
\mathbb B\times\mathbb U
\]
be a Mealy morphism.  Its terminal program category is
\[
\mathsf{Prog}_\Phi.
\]
It has input and output labellings
\[
I_\Phi:
\mathsf{Prog}_\Phi\to
(\mathbb A\times\mathbb U)^{\mathrm{op}},
\]
and
\[
\Omega_\Phi:
\mathsf{Prog}_\Phi\to
(\mathbb B\times\mathbb U)^{\mathrm{op}}.
\]
Write
\[
I_{\Phi,A},\ I_{\Phi,U},
\qquad
\Omega_{\Phi,B},\ \Omega_{\Phi,U}
\]
for their visible and hidden components.

\begin{definition}[Closed-loop program category]
\label{def:closed-loop-program-category}
The \emph{closed-loop program category} of \(\Phi\) over \(\mathbb U\) is
\[
\mathsf{Prog}^{\circ}_{\Phi}
:=
\operatorname{Eq}(I_{\Phi,U},\Omega_{\Phi,U}).
\]
Thus \(\mathsf{Prog}^{\circ}_{\Phi}\) is the internal subcategory of
\(\mathsf{Prog}_\Phi\) consisting of executions whose hidden input and hidden
output \(\mathbb U\)-labels agree.
\end{definition}

Equivalently, it is characterized by
\[
\begin{tikzcd}[column sep=large,row sep=large]
\mathsf{Prog}^{\circ}_{\Phi}
  \arrow[r, hook]
  \arrow[d]
&
\mathsf{Prog}_{\Phi}
  \arrow[d, "{(I_{\Phi,U},\Omega_{\Phi,U})}"]
\\
\mathbb U^{\mathrm{op}}
  \arrow[r, "\Delta_{\mathbb U^{\mathrm{op}}}"']
&
\mathbb U^{\mathrm{op}}\times\mathbb U^{\mathrm{op}} .
\end{tikzcd}
\]

\begin{proposition}[Dynamic trace-class criterion]
\label{prop:dynamic-trace-class-criterion}
For
\[
\Phi:
\mathbb A\times\mathbb U
\rightsquigarrow
\mathbb B\times\mathbb U,
\]
the following are equivalent:
\begin{enumerate}
\item \(\Phi\) is trace-class over \(\mathbb U\);
\item the restricted input projection
\[
I_{\Phi,A}^{\circ}:
\mathsf{Prog}^{\circ}_{\Phi}
\to
\mathbb A^{\mathrm{op}}
\]
is input-discrete;
\item the internal feedback equation
\[
h=\omega_U((a,h),x)
\]
has a unique solution for every admissible visible input and closed state,
equivalently
\[
W_\Phi\to D_\Phi
\]
is an isomorphism.
\end{enumerate}
\end{proposition}

\begin{proof}
The terminal program category of \(\Phi\) is precisely the labelled
input-discrete internal category representing \(\Phi\).  Equalizing the hidden
input and output labellings gives the closed-loop category
\[
\mathsf{Prog}^{\circ}_{\Phi}
=
\mathbb P^\circ.
\]
Thus the restricted input projection is input-discrete if and only if
\(\Phi\) is trace-class.  The equivalence with the unique feedback equation is
Proposition~\ref{prop:internal-trace-class-criterion}.
\end{proof}

\begin{theorem}[Program category of the traced Mealy morphism]
\label{thm:program-category-of-traced-mealy}
If
\[
\Phi:
\mathbb A\times\mathbb U
\rightsquigarrow
\mathbb B\times\mathbb U
\]
is trace-class over \(\mathbb U\), then there is a canonical isomorphism of
internal categories
\[
\mathsf{Prog}_{\operatorname{Tr}^{\mathbb U}(\Phi)}
\cong
\mathsf{Prog}^{\circ}_{\Phi}.
\]
Under this isomorphism, the input and output labellings of the traced program
category are the visible components
\[
I_{\Phi,A}^{\circ},
\qquad
\Omega_{\Phi,B}^{\circ}
\]
of the input and output labellings of \(\mathsf{Prog}_\Phi\).
\end{theorem}

\begin{proof}
By Definition~\ref{def:partial-trace-mealy},
\[
\operatorname{Tr}^{\mathbb U}(\Phi)
\]
is represented by the labelled input-discrete span
\[
\begin{tikzcd}[column sep=large]
\mathbb A^{\mathrm{op}}
&
\mathbb P^\circ
  \arrow[l, "I_A"']
  \arrow[r, "\Omega_B"]
&
\mathbb B^{\mathrm{op}} .
\end{tikzcd}
\]
The terminal program category associated to this Mealy morphism is therefore
\(\mathbb P^\circ\).  But \(\mathbb P^\circ\) is exactly
\[
\mathsf{Prog}^{\circ}_{\Phi}
=
\operatorname{Eq}(I_{\Phi,U},\Omega_{\Phi,U}).
\]
The statement about labellings follows from the definitions.
\end{proof}

Let
\[
\mathbf m_{\operatorname{Tr}^{\mathbb U}(\Phi)}
\]
be the primitive execution span of the traced Mealy morphism.  Let
\[
\mathbf m_{\Phi}^{\circ}
\]
be the primitive execution span of
\[
\mathsf{Prog}^{\circ}_{\Phi}.
\]

\begin{corollary}[Diamonds and boxes of feedback closure]
\label{cor:diamonds-boxes-feedback-closure}
Assume the predicate-transformer structure of
Chapter~\ref{ch:relative-mealy-program-categories} is available.

If \(\Phi\) is trace-class over \(\mathbb U\), then for every predicate
\[
\varphi\hookrightarrow P^\circ
\]
one has
\[
\langle
\mathbf m_{\operatorname{Tr}^{\mathbb U}(\Phi)}
\rangle\varphi
=
\langle
\mathbf m_{\Phi}^{\circ}
\rangle\varphi.
\]
In the first-order Heyting setting one also has
\[
[
\mathbf m_{\operatorname{Tr}^{\mathbb U}(\Phi)}
]\varphi
=
[
\mathbf m_{\Phi}^{\circ}
]\varphi.
\]
\end{corollary}

\begin{proof}
By Theorem~\ref{thm:program-category-of-traced-mealy}, the primitive execution
span of the traced Mealy morphism is canonically isomorphic to the primitive
execution span of the closed-loop program category.  The diamond and box
predicate transformers are invariant under isomorphism of spans.
\end{proof}

\begin{remark}[Operational reading]
\label{rem:operational-reading-feedback-closure}
In generalized-element notation,
\[
x\models
\langle
\mathbf m_{\operatorname{Tr}^{\mathbb U}(\Phi)}
\rangle\varphi
\]
means that there exists a visible input arrow \(a\) admissible at \(x\) such
that, writing
\[
h_\Phi(a,x)
\]
for the unique feedback solution,
\[
\delta((a,h_\Phi(a,x)),x)\models\varphi.
\]
The uniqueness of \(h_\Phi(a,x)\) is the trace-class condition that makes the
closed-loop span Mealy rather than merely relational.
\end{remark}

\section{Relative semantics of traced systems}
\label{sec:relative-semantics-traced-systems}

Let
\[
Y=(Y\xrightarrow{r}B_0,\sigma)
\]
be a right \(\mathbb B\)-action.  If
\[
\Phi:
\mathbb A\times\mathbb U
\rightsquigarrow
\mathbb B\times\mathbb U
\]
is trace-class over \(\mathbb U\), then
\[
\operatorname{Tr}^{\mathbb U}(\Phi):
\mathbb A\rightsquigarrow\mathbb B
\]
has the relative program category
\[
\mathsf{Prog}_{\operatorname{Tr}^{\mathbb U}(\Phi),Y}.
\]

On the other hand, the untraced Mealy morphism \(\Phi\) may be interpreted
relative to the downstream \((\mathbb B\times\mathbb U)\)-action
\[
Y\times\mathbf 1_{\mathbb U}.
\]
The terminal right \(\mathbb U\)-action is
\[
\mathbf 1_{\mathbb U}
=
(U_0\xrightarrow{1_{U_0}}U_0,\tau_U),
\]
with
\[
\tau_U(h,t_Uh)=s_Uh.
\]
Thus
\[
Y\times\mathbf 1_{\mathbb U}
\]
has state object \(Y\times U_0\), structure map
\[
Y\times U_0\to B_0\times U_0,
\qquad
(y,u)\longmapsto(r(y),u),
\]
and action
\[
(\beta,h),(y,t_Uh)
\longmapsto
(\sigma(\beta,y),s_Uh).
\]

The relative program category of the untraced system is
\[
\mathsf{Prog}_{\Phi,Y\times\mathbf 1_{\mathbb U}}.
\]
It has an input labelling
\[
I_{\Phi,Y\times\mathbf 1_{\mathbb U}}:
\mathsf{Prog}_{\Phi,Y\times\mathbf 1_{\mathbb U}}
\to
(\mathbb A\times\mathbb U)^{\mathrm{op}},
\]
and a downstream comparison functor
\[
\Lambda_{\Phi,Y\times\mathbf 1_{\mathbb U}}:
\mathsf{Prog}_{\Phi,Y\times\mathbf 1_{\mathbb U}}
\to
\int_{\mathbb B\times\mathbb U}(Y\times\mathbf 1_{\mathbb U}).
\]
The hidden input labelling is the \(\mathbb U\)-component of
\(I_{\Phi,Y\times\mathbf 1_{\mathbb U}}\).  The hidden output labelling is the
composite
\[
\mathsf{Prog}_{\Phi,Y\times\mathbf 1_{\mathbb U}}
\xrightarrow{\Lambda_{\Phi,Y\times\mathbf 1_{\mathbb U}}}
\int_{\mathbb B\times\mathbb U}(Y\times\mathbf 1_{\mathbb U})
\longrightarrow
\mathbb U^{\mathrm{op}},
\]
where the last functor projects the downstream action category to its terminal
\(\mathbb U\)-action component.

\begin{definition}[Relative closed-loop program category]
\label{def:relative-closed-loop-program-category}
The relative closed-loop program category
\[
\left(
\mathsf{Prog}_{\Phi,Y\times\mathbf 1_{\mathbb U}}
\right)^\circ
\]
is the equalizer of the hidden input and hidden output labellings defined
above.
\end{definition}

Inside it, feedback closure over \(\mathbb U\) imposes
\[
p_U=q_U
\]
on objects and
\[
h=\omega_U((a,h),x)
\]
on arrows.

\begin{theorem}[Relative semantics commutes with feedback closure]
\label{thm:relative-semantics-commutes-feedback}
If \(\Phi\) is trace-class over \(\mathbb U\), then there is a canonical
isomorphism
\[
\mathsf{Prog}_{\operatorname{Tr}^{\mathbb U}(\Phi),Y}
\cong
\left(
\mathsf{Prog}_{\Phi,Y\times\mathbf 1_{\mathbb U}}
\right)^\circ.
\]
\end{theorem}

\begin{proof}
The state object of the left-hand side is
\[
P^\circ\times_{q_B,B_0,r}Y.
\]
The state object of the right-hand side before feedback closure is
\[
P\times_{B_0\times U_0}(Y\times U_0),
\]
where the map
\[
P\to B_0\times U_0
\]
is \((q_B,q_U)\), and the map
\[
Y\times U_0\to B_0\times U_0
\]
is \((r,1_{U_0})\).  Hence this state object is canonically isomorphic to
\[
P\times_{q_B,B_0,r}Y.
\]
Under this identification, the hidden terminal-action coordinate is \(q_U(x)\).
Feedback closure imposes
\[
p_U(x)=q_U(x).
\]
Thus the closed-loop state object is
\[
P^\circ\times_{q_B,B_0,r}Y.
\]

On arrows, the left-hand side uses the unique feedback solution
\[
h_\Phi(a,x)
\]
to process a visible input arrow \(a\).  The right-hand side consists
precisely of the closed-loop executions of the untraced relative program
category.  The downstream terminal \(\mathbb U\)-action updates the hidden
downstream coordinate from \(t_Uh\) to \(s_Uh\), and the hidden output
labelling records \(\omega_U((a,h),x)\).  Hence the equalizer condition is
exactly
\[
h=\omega_U((a,h),x).
\]
Trace-class says that such a hidden feedback arrow exists uniquely.  Therefore
the primitive execution objects, source maps, target maps, identities, and
compositions agree under the canonical identification.
\end{proof}

\section{Polynomial response coalgebras and feedback trace}
\label{sec:polynomial-response-form-feedback}

We record compatibility with the polynomial transpose constructed in
Chapter~\ref{ch:polynomial-relative-mealy-actions}.  In this section, the
polynomial hypotheses of that chapter are assumed.

The point is not merely that the traced response profile can be written by
substitution.  There is a partial trace operation on multiplicative response
coalgebras, and the Mealy-induced coalgebras are stable under it.

Let
\[
\mathbb A,\mathbb B,\mathbb U
\]
be internal categories.  A multiplicative response coalgebra of hidden type
\[
\mathbb A\times\mathbb U
\longrightarrow
\mathbb B\times\mathbb U
\]
consists of a state span
\[
A_0\times U_0
\xleftarrow{p}
S
\xrightarrow{q}
B_0\times U_0,
\]
written
\[
p=(p_A,p_U),
\qquad
q=(q_B,q_U),
\]
together with a response operation assigning to every admissible triple
\[
(a,h,x),
\qquad
t_A(a)=p_A(x),
\qquad
t_U(h)=p_U(x),
\]
a next state
\[
\partial(a,h,x)\in S
\]
and output arrows
\[
\beta(a,h,x)\in B_1,
\qquad
\kappa(a,h,x)\in U_1,
\]
such that
\[
p_A\partial(a,h,x)=s_A(a),
\qquad
p_U\partial(a,h,x)=s_U(h),
\]
\[
s_B\beta(a,h,x)=q_B\partial(a,h,x),
\qquad
t_B\beta(a,h,x)=q_B(x),
\]
and
\[
s_U\kappa(a,h,x)=q_U\partial(a,h,x),
\qquad
t_U\kappa(a,h,x)=q_U(x).
\]
The unit and multiplication laws are the curried forms of the Mealy unit and
multiplication laws:
\[
\partial(e_Ap_Ax,e_Up_Ux,x)=x,
\]
\[
\beta(e_Ap_Ax,e_Up_Ux,x)=e_Bq_Bx,
\qquad
\kappa(e_Ap_Ax,e_Up_Ux,x)=e_Uq_Ux,
\]
and, for composable visible and hidden inputs
\[
a:u\to v,\quad b:v\to w,
\qquad
h:u'\to v',\quad k:v'\to w',
\]
with
\[
p_A(x)=w,
\qquad
p_U(x)=w',
\]
one has
\[
\partial(b\circ a,k\circ h,x)
=
\partial(a,h,\partial(b,k,x)),
\]
\[
\beta(b\circ a,k\circ h,x)
=
\beta(b,k,x)\circ
\beta(a,h,\partial(b,k,x)),
\]
and
\[
\kappa(b\circ a,k\circ h,x)
=
\kappa(b,k,x)\circ
\kappa(a,h,\partial(b,k,x)).
\]

\begin{definition}[Trace class for response coalgebras]
\label{def:trace-class-response-coalgebra}
Let
\[
c=(S,p,q,\partial,\beta,\kappa)
\]
be a multiplicative response coalgebra of hidden type
\[
\mathbb A\times\mathbb U
\longrightarrow
\mathbb B\times\mathbb U.
\]
Define the closed state object
\[
S^\circ:=\operatorname{Eq}(p_U,q_U).
\]
Let
\[
D_c:=A_1\times_{t_A,A_0,p_A}S^\circ,
\]
and
\[
H_c:=D_c\times_{p_U\pi_S,U_0,t_U}U_1.
\]
There are two maps
\[
H_c\to U_1.
\]
The first is the hidden-input projection
\[
(a,x,h)\longmapsto h.
\]
The second is the hidden response
\[
(a,x,h)\longmapsto \kappa(a,h,x).
\]
The coalgebra \(c\) is \emph{trace-class over \(\mathbb U\)} if the equalizer
\[
W_c
:=
\operatorname{Eq}\bigl(h,\kappa(a,h,x)\bigr)
\]
projects isomorphically onto \(D_c\).
\end{definition}

When \(c\) is trace-class, write
\[
h_c(a,x)
\]
for the unique hidden feedback solution.

\begin{definition}[Partial trace of response coalgebras]
\label{def:partial-trace-response-coalgebras}
Let
\[
c=(S,p,q,\partial,\beta,\kappa)
\]
be a trace-class multiplicative response coalgebra of hidden type
\[
\mathbb A\times\mathbb U
\longrightarrow
\mathbb B\times\mathbb U.
\]
Its partial trace
\[
\operatorname{Tr}^{\mathbb U}(c)
\]
is the multiplicative response coalgebra of type
\[
\mathbb A\longrightarrow\mathbb B
\]
with state span
\[
A_0
\xleftarrow{p_A}
S^\circ
\xrightarrow{q_B}
B_0,
\]
next-state map
\[
\partial^{\operatorname{tr}}(a,x)
=
\partial(a,h_c(a,x),x),
\]
and output map
\[
\beta^{\operatorname{tr}}(a,x)
=
\beta(a,h_c(a,x),x).
\]
\end{definition}

\begin{proposition}[Well-definedness of coalgebra trace]
\label{prop:coalgebra-trace-well-defined}
If \(c\) is trace-class over \(\mathbb U\), then
\[
\operatorname{Tr}^{\mathbb U}(c)
\]
is a multiplicative response coalgebra of type
\[
\mathbb A\longrightarrow\mathbb B.
\]
\end{proposition}

\begin{proof}
Let
\[
h=h_c(a,x).
\]
The typing equations for \(c\) give
\[
p_A\partial(a,h,x)=s_A(a),
\qquad
p_U\partial(a,h,x)=s_U(h),
\]
and
\[
q_U\partial(a,h,x)=s_U\kappa(a,h,x).
\]
Since \(h=\kappa(a,h,x)\), the two hidden state objects agree:
\[
p_U\partial(a,h,x)
=
q_U\partial(a,h,x).
\]
Thus
\[
\partial^{\operatorname{tr}}(a,x)\in S^\circ.
\]
The visible output typing equations are inherited from those of \(c\):
\[
s_B\beta^{\operatorname{tr}}(a,x)
=
q_B\partial^{\operatorname{tr}}(a,x),
\qquad
t_B\beta^{\operatorname{tr}}(a,x)
=
q_B(x).
\]

For the unit law, take a closed state \(x\in S^\circ\).  The hidden identity
\[
e_Up_Ux
\]
solves the feedback equation for the visible identity input, because the unit
law for \(c\) gives
\[
\kappa(e_Ap_Ax,e_Up_Ux,x)=e_Uq_Ux=e_Up_Ux.
\]
By trace-class uniqueness,
\[
h_c(e_Ap_Ax,x)=e_Up_Ux.
\]
Substituting this into the unit laws of \(c\) gives the unit laws for
\[
\operatorname{Tr}^{\mathbb U}(c).
\]

For multiplication, let
\[
a:u\to v,
\qquad
b:v\to w,
\qquad
p_A(x)=w.
\]
Write
\[
h_b:=h_c(b,x),
\qquad
x_b:=\partial(b,h_b,x),
\qquad
h_a:=h_c(a,x_b).
\]
Then
\[
t_U(h_b)=p_U(x),
\qquad
t_U(h_a)=p_U(x_b)=s_U(h_b),
\]
so the composite
\[
h_b\circ h_a
\]
is an admissible hidden input for the visible composite \(b\circ a\) at \(x\).
By the multiplication law for \(c\),
\[
\kappa(b\circ a,h_b\circ h_a,x)
=
\kappa(b,h_b,x)\circ
\kappa(a,h_a,x_b).
\]
Since \(h_b\) and \(h_a\) are feedback solutions, this becomes
\[
\kappa(b\circ a,h_b\circ h_a,x)
=
h_b\circ h_a.
\]
Thus \(h_b\circ h_a\) is a feedback solution for the composite input.  By
trace-class uniqueness,
\[
h_c(b\circ a,x)=h_b\circ h_a.
\]
Substituting this identity into the multiplication laws for
\(\partial\) and \(\beta\) gives
\[
\partial^{\operatorname{tr}}(b\circ a,x)
=
\partial^{\operatorname{tr}}(a,\partial^{\operatorname{tr}}(b,x)),
\]
and
\[
\beta^{\operatorname{tr}}(b\circ a,x)
=
\beta^{\operatorname{tr}}(b,x)\circ
\beta^{\operatorname{tr}}(a,\partial^{\operatorname{tr}}(b,x)).
\]
Hence the traced structure is multiplicative.
\end{proof}

\begin{proposition}[Trace of Mealy-induced response coalgebras]
\label{prop:trace-mealy-induced-coalgebras}
Let
\[
\Phi=(P,\delta,\omega):
\mathbb A\times\mathbb U
\rightsquigarrow
\mathbb B\times\mathbb U
\]
be a Mealy morphism, and let
\[
c_\Phi
\]
be its induced multiplicative response coalgebra from
Chapter~\ref{ch:polynomial-relative-mealy-actions}.  Then \(\Phi\) is
trace-class over \(\mathbb U\) if and only if \(c_\Phi\) is trace-class over
\(\mathbb U\).  When this holds,
\[
\operatorname{Tr}^{\mathbb U}(c_\Phi)
=
c_{\operatorname{Tr}^{\mathbb U}(\Phi)}.
\]
\end{proposition}

\begin{proof}
For the Mealy-induced response coalgebra, the response data are
\[
\partial(a,h,x)=\delta((a,h),x),
\]
\[
\beta(a,h,x)=\omega_B((a,h),x),
\qquad
\kappa(a,h,x)=\omega_U((a,h),x).
\]
Therefore the coalgebra feedback equation
\[
h=\kappa(a,h,x)
\]
is exactly the Mealy feedback equation
\[
h=\omega_U((a,h),x).
\]
The coalgebra closed state object is
\[
\operatorname{Eq}(p_U,q_U),
\]
which is the closed state object of \(\Phi\).  Hence the trace-class
conditions agree.

When the trace-class condition holds, both traced structures have carrier
\[
P^\circ=\operatorname{Eq}(p_U,q_U).
\]
Their next-state and visible-output maps are, respectively,
\[
\delta((a,h_\Phi(a,x)),x)
\]
and
\[
\omega_B((a,h_\Phi(a,x)),x).
\]
Thus the traced coalgebra is exactly the coalgebra induced by the traced Mealy
morphism.
\end{proof}

\begin{theorem}[Flattening commutes with feedback trace]
\label{thm:flattening-commutes-feedback}
Let
\[
\Phi:
\mathbb A\times\mathbb U
\rightsquigarrow
\mathbb B\times\mathbb U
\]
be trace-class over \(\mathbb U\).  Then
\[
\operatorname{Flat}
\left(
\operatorname{Tr}^{\mathbb U}(c_\Phi)
\right)
\cong
\mathbf m_{\Phi}^{\circ}.
\]
Equivalently,
\[
\operatorname{Flat}
\left(c_{\operatorname{Tr}^{\mathbb U}(\Phi)}\right)
\cong
\mathbf m_{\operatorname{Tr}^{\mathbb U}(\Phi)}
\cong
\mathbf m_{\Phi}^{\circ}.
\]
\end{theorem}

\begin{proof}
By Proposition~\ref{prop:trace-mealy-induced-coalgebras},
\[
\operatorname{Tr}^{\mathbb U}(c_\Phi)
=
c_{\operatorname{Tr}^{\mathbb U}(\Phi)}.
\]
By Chapter~\ref{ch:polynomial-relative-mealy-actions}, flattening a
Mealy-induced response coalgebra recovers the primitive execution span of its
program category.  Therefore
\[
\operatorname{Flat}
\left(c_{\operatorname{Tr}^{\mathbb U}(\Phi)}\right)
\cong
\mathbf m_{\operatorname{Tr}^{\mathbb U}(\Phi)}.
\]
By Theorem~\ref{thm:program-category-of-traced-mealy}, the program category of
the traced Mealy morphism is the closed-loop program category of \(\Phi\).
Hence
\[
\mathbf m_{\operatorname{Tr}^{\mathbb U}(\Phi)}
\cong
\mathbf m_{\Phi}^{\circ}.
\]
Combining these canonical isomorphisms gives the result.
\end{proof}

The compatibility is expressed by the commutative diagram, up to the canonical
isomorphisms described above:
\[
\begin{tikzcd}[column sep=huge,row sep=large]
c_{\Phi}
  \arrow[r, maps to, "{\operatorname{Flat}}"]
  \arrow[d, maps to, "{\operatorname{Tr}^{\mathbb U}}"']
&
\mathbf m_{\Phi}
  \arrow[d, maps to, "{(-)^\circ}"]
\\
\operatorname{Tr}^{\mathbb U}(c_\Phi)
=
c_{\operatorname{Tr}^{\mathbb U}(\Phi)}
  \arrow[r, maps to, "{\operatorname{Flat}}"']
&
\mathbf m_{\Phi}^{\circ}.
\end{tikzcd}
\]

\section{A partial Geometry-of-Interaction paracategory}
\label{sec:partial-geometry-of-interaction-category}

The partial trace on Mealy morphisms gives a bidirectional process structure
by the partial analogue of the Joyal--Street--Verity Int-construction.  The
partial Int-construction and its paracategorical nature are developed by
Malherbe--Scott--Selinger using Freyd paracategories
\cite{MalherbeScottSelingerPartiallyTracedCategories}.  We now spell out the
resulting construction in the present Mealy setting.

\begin{definition}[Bidirectional Mealy interface]
\label{def:bidirectional-mealy-interface}
A \emph{bidirectional Mealy interface} is a pair
\[
\mathbb A=(\mathbb A^+,\mathbb A^-)
\]
of internal categories.

The dual interface is
\[
\mathbb A^\bot:=(\mathbb A^-,\mathbb A^+).
\]
The tensor product is componentwise:
\[
\mathbb A\otimes\mathbb B
:=
(\mathbb A^+\times\mathbb B^+,\,
 \mathbb A^-\times\mathbb B^-).
\]
The tensor unit is
\[
\mathbb I:=(\mathbf 1,\mathbf 1),
\]
where \(\mathbf 1\) is the terminal internal category.
\end{definition}

\begin{definition}[Bidirectional Mealy process]
\label{def:bidirectional-mealy-process}
A \emph{bidirectional Mealy process}
\[
F:\mathbb A\to\mathbb B
\]
is a Mealy morphism
\[
F:
\mathbb A^+\times\mathbb B^-
\rightsquigarrow
\mathbb B^+\times\mathbb A^-.
\]
Thus a process consumes a positive \(\mathbb A\)-input and a negative
\(\mathbb B\)-context, and emits a positive \(\mathbb B\)-output and a
negative \(\mathbb A\)-response.
\end{definition}

\begin{definition}[Partial Int composition]
\label{def:partial-int-composition-mealy}
Let
\[
F:\mathbb A\to\mathbb B
\]
and
\[
G:\mathbb B\to\mathbb C
\]
be bidirectional Mealy processes, so
\[
F:
\mathbb A^+\times\mathbb B^-
\rightsquigarrow
\mathbb B^+\times\mathbb A^-,
\]
and
\[
G:
\mathbb B^+\times\mathbb C^-
\rightsquigarrow
\mathbb C^+\times\mathbb B^-.
\]
Their \emph{open pipeline} is the Mealy morphism
\[
H_{G,F}:
\mathbb A^+\times\mathbb C^-\times\mathbb B^-
\rightsquigarrow
\mathbb C^+\times\mathbb A^-\times\mathbb B^-
\]
defined, up to the canonical associativity and symmetry constraints of the
symmetric monoidal category \(\mathsf{Mly}(\mathcal E)\), by
\[
H_{G,F}
:=
(1_{\mathbb C^+}\times\sigma_{\mathbb B^-,\mathbb A^-})
\circ
(G\times 1_{\mathbb A^-})
\circ
(1_{\mathbb B^+}\times\sigma_{\mathbb A^-,\mathbb C^-})
\circ
(F\times 1_{\mathbb C^-})
\circ
(1_{\mathbb A^+}\times\sigma_{\mathbb C^-,\mathbb B^-}).
\]
The forward \(\mathbb B^+\)-wire is connected by ordinary Mealy composition
inside the pipeline, and the remaining open middle interface is the
counterflow wire \(\mathbb B^-\).

The partial composite
\[
G\circ_{\partial}F
\]
is defined when
\[
H_{G,F}
\]
is trace-class over \(\mathbb B^-\).  In that case,
\[
G\circ_{\partial}F
:=
\operatorname{Tr}^{\mathbb B^-}(H_{G,F}).
\]
\end{definition}

The structural rebracketing underlying \(H_{G,F}\) is represented by
\[
\begin{tikzcd}[column sep=huge,row sep=large]
\mathbb A^+\times\mathbb C^-\times\mathbb B^-
  \arrow[r, "{1\times\sigma}"]
&
\mathbb A^+\times\mathbb B^-\times\mathbb C^-
  \arrow[r, "{F\times 1}"]
&
\mathbb B^+\times\mathbb A^-\times\mathbb C^-
  \arrow[r, "{1\times\sigma}"]
&
\mathbb B^+\times\mathbb C^-\times\mathbb A^-
  \arrow[r, "{G\times 1}"]
&
\mathbb C^+\times\mathbb B^-\times\mathbb A^-
  \arrow[r, "{1\times\sigma}"]
&
\mathbb C^+\times\mathbb A^-\times\mathbb B^- .
\end{tikzcd}
\]
The trace over \(\mathbb B^-\) then has type
\[
\operatorname{Tr}^{\mathbb B^-}(H_{G,F}):
\mathbb A^+\times\mathbb C^-
\rightsquigarrow
\mathbb C^+\times\mathbb A^-.
\]

\begin{definition}[Partial Int paracategory]
\label{def:partial-int-paracategory-mealy}
The \emph{partial Int construction} of the Mealy category is the paracategory
\[
\operatorname{Int}_{\partial}(\mathsf{Mly}(\mathcal E))
\]
whose objects are bidirectional Mealy interfaces, whose arrows are
bidirectional Mealy processes, and whose partial composition is
Definition~\ref{def:partial-int-composition-mealy}.
\end{definition}

\begin{theorem}[Partial Geometry of Interaction for Mealy morphisms]
\label{thm:partial-goi-mealy}
The structure
\[
\operatorname{Int}_{\partial}(\mathsf{Mly}(\mathcal E))
\]
is a symmetric monoidal paracategory.  Its partial composition is obtained by
forming the open pipeline and then partially tracing the shared counterflow
interface.
\end{theorem}

\begin{proof}
The construction is the partial Int-construction applied to the partially
traced symmetric monoidal category
\[
\mathsf{Mly}(\mathcal E)
\]
of Theorem~\ref{thm:partial-trace-on-mealy}.  The paracategory laws follow
from the partial trace axioms, exactly as in the partial Int-construction of
Malherbe--Scott--Selinger
\cite{MalherbeScottSelingerPartiallyTracedCategories}.  Symmetric monoidal
structure is componentwise:
\[
(\mathbb A^+,\mathbb A^-)\otimes(\mathbb B^+,\mathbb B^-)
=
(\mathbb A^+\times\mathbb B^+,\mathbb A^-\times\mathbb B^-).
\]
Associativity of partial composition holds whenever both sides are defined,
because the two composites are obtained by tracing canonically isomorphic
open pipelines over canonically isomorphic hidden interfaces.  This is the
vanishing and associativity content of the inherited partial trace.
\end{proof}

\section{Multiplicative partial linear logic}
\label{sec:multiplicative-partial-linear-logic-mealy}

We now give a natural-deduction presentation of the multiplicative partial
linear logic interpreted by
\[
\operatorname{Int}_{\partial}(\mathsf{Mly}(\mathcal E)).
\]
The calculus is linear: assumptions are used exactly once, exchange is allowed,
and cut is a partial rule governed by trace-class feedback.

\subsection{Formulas, contexts, and interpretations}

Fix a set of atomic formulas.  Each atom \(X\) is assigned a bidirectional
Mealy interface
\[
\llbracket X\rrbracket=(\mathbb X^+,\mathbb X^-).
\]
Formulas are generated by
\[
A,B ::= X \mid A^\bot \mid A\otimes B \mid \mathbf 1
        \mid A\parr B \mid \bot .
\]
Linear negation is involutive and satisfies the De Morgan equations
\[
(A^\bot)^\bot=A,
\]
\[
(A\otimes B)^\bot=A^\bot\parr B^\bot,
\qquad
(A\parr B)^\bot=A^\bot\otimes B^\bot,
\]
\[
\mathbf 1^\bot=\bot,
\qquad
\bot^\bot=\mathbf 1.
\]
The linear implication abbreviation is
\[
A\multimap B:=A^\bot\parr B.
\]

The interpretation of formulas is defined by
\[
\llbracket A^\bot\rrbracket
:=
\llbracket A\rrbracket^\bot,
\]
\[
\llbracket A\otimes B\rrbracket
:=
\llbracket A\rrbracket\otimes\llbracket B\rrbracket,
\]
\[
\llbracket A\parr B\rrbracket
:=
\left(
\llbracket A^\bot\rrbracket
\otimes
\llbracket B^\bot\rrbracket
\right)^\bot,
\]
and
\[
\llbracket \mathbf 1\rrbracket
=
\llbracket \bot\rrbracket
=
\mathbb I.
\]
Accordingly,
\[
\llbracket A\multimap B\rrbracket
=
\llbracket A^\bot\parr B\rrbracket.
\]

A context is a finite list of distinct typed variables
\[
\Gamma=x_1:A_1,\ldots,x_n:A_n.
\]
Its interpretation is
\[
\llbracket \Gamma\rrbracket
:=
\llbracket A_1\rrbracket\otimes\cdots\otimes\llbracket A_n\rrbracket,
\]
with the empty context interpreted as
\[
\llbracket \cdot\rrbracket=\mathbb I.
\]
Exchange is interpreted by the symmetry isomorphisms of
\[
\operatorname{Int}_{\partial}(\mathsf{Mly}(\mathcal E)).
\]

A judgement
\[
\Gamma\vdash M:A
\]
is interpreted as a bidirectional Mealy process
\[
\llbracket M\rrbracket:
\llbracket \Gamma\rrbracket
\longrightarrow
\llbracket A\rrbracket.
\]
Thus a proof term in context \(\Gamma\) denotes a process from the tensor of
the input interfaces in \(\Gamma\) to the output interface \(A\).

\subsection{Structural transposition and evaluation}

For bidirectional interfaces \(\mathbb X,\mathbb A,\mathbb B\), the partial
Int construction gives a canonical structural bijection
\[
\operatorname{Int}_{\partial}(\mathsf{Mly})(\mathbb X\otimes\mathbb A,\mathbb B)
\cong
\operatorname{Int}_{\partial}(\mathsf{Mly})(\mathbb X,\mathbb A\multimap\mathbb B),
\]
where
\[
\mathbb A\multimap\mathbb B:=\mathbb A^\bot\parr\mathbb B.
\]
This bijection is implemented entirely by associativity, unit, symmetry, and
polarity-rebracketing isomorphisms of product interfaces.  Hence it is total;
no feedback equation is solved in forming the transpose.

Explicitly, a process
\[
F:\mathbb X\otimes\mathbb A\to\mathbb B
\]
is a Mealy morphism
\[
(\mathbb X^+\times\mathbb A^+)\times\mathbb B^-
\rightsquigarrow
\mathbb B^+\times(\mathbb X^-\times\mathbb A^-).
\]
After rebracketing and permuting factors, this is equivalently a Mealy
morphism
\[
\mathbb X^+\times(\mathbb A^+\times\mathbb B^-)
\rightsquigarrow
(\mathbb A^-\times\mathbb B^+)\times\mathbb X^-,
\]
which is precisely a process
\[
\mathbb X\to\mathbb A\multimap\mathbb B.
\]

The evaluation process
\[
\operatorname{ev}_{A,B}:
(A\multimap B)\otimes A\to B
\]
is the stateless bidirectional Mealy process induced by the canonical product
permutation
\[
(A^-\times B^+)\times A^+\times B^-
\cong
B^+\times A^+\times B^-\times A^-.
\]
Equivalently, as a Mealy morphism, it is the representable Mealy morphism
carried by this permutation of positive and negative interface factors.  Its
transition is forced by the source projection and its output is the
corresponding functorial transport along the product symmetry.

\subsection{Trace-class substitution}

The primitive partial operation of the proof system is linear substitution.
Let
\[
\Gamma\vdash M:A
\]
and
\[
\Delta,x:A\vdash N:B
\]
be derivations.  The semantic open pipeline for substituting \(M\) for \(x\)
in \(N\) is
\[
\llbracket \Delta\rrbracket\otimes\llbracket \Gamma\rrbracket
\xrightarrow{1\otimes \llbracket M\rrbracket}
\llbracket \Delta\rrbracket\otimes\llbracket A\rrbracket
\xrightarrow{\llbracket N\rrbracket}
\llbracket B\rrbracket,
\]
read inside the partial Int paracategory.  The substitution
\[
N[M/x]
\]
is defined precisely when this open pipeline is trace-class along the cut
interface \(\llbracket A\rrbracket\).  When defined, its denotation is
\[
\llbracket N[M/x]\rrbracket
=
\llbracket N\rrbracket\circ_{\partial}
(1_{\llbracket\Delta\rrbracket}\otimes \llbracket M\rrbracket).
\]
We write this side condition as
\[
\operatorname{TC}_x(N,M).
\]

\begin{definition}[Partial cut/substitution rule]
\label{def:natural-deduction-partial-cut}
The partial cut rule is
\[
\frac{
  \Gamma\vdash M:A
  \qquad
  \Delta,x:A\vdash N:B
  \qquad
  \operatorname{TC}_x(N,M)
}{
  \Delta,\Gamma\vdash N[M/x]:B
}
\quad(\operatorname{cut}).
\]
\end{definition}

All elimination rules below are instances of this partial substitution rule
applied to canonical elimination processes.

\subsection{Identity, exchange, and structural linearity}

The identity rule is
\[
\frac{}{x:A\vdash x:A}
\quad(\operatorname{id}).
\]
Its denotation is the identity process
\[
1_{\llbracket A\rrbracket}:
\llbracket A\rrbracket\to\llbracket A\rrbracket.
\]

The exchange rule is
\[
\frac{
  x_1:A_1,\ldots,x_i:A_i,x_{i+1}:A_{i+1},\ldots,x_n:A_n
  \vdash M:B
}{
  x_1:A_1,\ldots,x_{i+1}:A_{i+1},x_i:A_i,\ldots,x_n:A_n
  \vdash \sigma_i M:B
}
\quad(\operatorname{exch}),
\]
interpreted by precomposition with the corresponding symmetry isomorphism.
There are no weakening or contraction rules.

\subsection{Tensor rules}

Tensor introduction is
\[
\frac{
  \Gamma\vdash M:A
  \qquad
  \Delta\vdash N:B
}{
  \Gamma,\Delta\vdash M\otimes N:A\otimes B
}
\quad(\otimes I),
\]
with denotation
\[
\llbracket M\otimes N\rrbracket
=
\llbracket M\rrbracket\otimes\llbracket N\rrbracket.
\]

Tensor elimination is
\[
\frac{
  \Gamma\vdash M:A\otimes B
  \qquad
  \Delta,x:A,y:B\vdash N:C
  \qquad
  \operatorname{TC}_{x,y}(N,M)
}{
  \Delta,\Gamma\vdash
  \operatorname{let}\ x\otimes y=M\ \operatorname{in}\ N
  :C
}
\quad(\otimes E).
\]
Here \(\operatorname{TC}_{x,y}(N,M)\) means that the open pipeline obtained by
cutting the \(A\otimes B\)-output of \(M\) against the \(x:A,y:B\)-input of
\(N\) is trace-class.  Equivalently, after the canonical associativity
identification
\[
\llbracket A\otimes B\rrbracket
=
\llbracket A\rrbracket\otimes\llbracket B\rrbracket,
\]
this is the partial substitution condition for the pair \(x,y\).

The beta equation for tensor is
\[
\operatorname{let}\ x\otimes y=(M\otimes N)\ \operatorname{in}\ L
=
L[M/x,N/y],
\]
whenever the indicated substitutions are trace-class.  The eta equation is
\[
\operatorname{let}\ x\otimes y=P\ \operatorname{in}\ x\otimes y
=
P.
\]
Both equations are interpreted by the tensor-unit and associativity axioms of
the partial Int paracategory.

\subsection{Unit rules}

The unit introduction rule is
\[
\frac{}{\cdot\vdash \star:\mathbf 1}
\quad(\mathbf 1 I),
\]
interpreted by the identity process
\[
1_{\mathbb I}:\mathbb I\to\mathbb I.
\]

The unit elimination rule is
\[
\frac{
  \Gamma\vdash M:\mathbf 1
  \qquad
  \Delta\vdash N:C
  \qquad
  \operatorname{TC}_{\mathbf 1}(N,M)
}{
  \Delta,\Gamma\vdash
  \operatorname{let}\ \star=M\ \operatorname{in}\ N
  :C
}
\quad(\mathbf 1 E).
\]
The beta and eta equations are
\[
\operatorname{let}\ \star=\star\ \operatorname{in}\ N=N,
\]
and
\[
\operatorname{let}\ \star=M\ \operatorname{in}\ \star=M,
\]
whenever the displayed cuts are defined.

The par unit is
\[
\bot=\mathbf 1^\bot.
\]
Its canonical closed proof is
\[
\frac{}{\cdot\vdash \star^\bot:\bot}
\quad(\bot I),
\]
interpreted by the same unit interface \(\mathbb I\), with opposite polarity.

\subsection{Linear implication and par rules}

The implication connective is the abbreviation
\[
A\multimap B=A^\bot\parr B.
\]
Its introduction rule is currying:
\[
\frac{
  \Gamma,x:A\vdash M:B
}{
  \Gamma\vdash \lambda x.M:A\multimap B
}
\quad(\multimap I).
\]
Semantically, if
\[
\llbracket M\rrbracket:
\llbracket\Gamma\rrbracket\otimes\llbracket A\rrbracket
\to
\llbracket B\rrbracket,
\]
then
\[
\llbracket \lambda x.M\rrbracket:
\llbracket\Gamma\rrbracket
\to
\llbracket A\multimap B\rrbracket
\]
is its structural Int-transpose.

Implication elimination is application:
\[
\frac{
  \Gamma\vdash F:A\multimap B
  \qquad
  \Delta\vdash M:A
  \qquad
  \operatorname{TC}_{\operatorname{ev}}(F,M)
}{
  \Gamma,\Delta\vdash F\,M:B
}
\quad(\multimap E).
\]
The side condition says that the evaluation pipeline
\[
\llbracket \Gamma\rrbracket\otimes\llbracket\Delta\rrbracket
\xrightarrow{\llbracket F\rrbracket\otimes\llbracket M\rrbracket}
\llbracket A\multimap B\rrbracket\otimes\llbracket A\rrbracket
\xrightarrow{\operatorname{ev}_{A,B}}
\llbracket B\rrbracket
\]
is trace-class.  The beta equation is
\[
(\lambda x.M)\,N=M[N/x],
\]
whenever the right-hand substitution is trace-class.  The eta equation is
\[
\lambda x.(F\,x)=F,
\]
whenever the displayed evaluation is trace-class.

Since
\[
A\parr B=A^\bot\multimap B,
\]
the par rules are the corresponding implication rules specialized to the
negative antecedent.

Par introduction is
\[
\frac{
  \Gamma,x:A^\bot\vdash M:B
}{
  \Gamma\vdash \lambda x.M:A\parr B
}
\quad(\parr I).
\]
Par elimination is
\[
\frac{
  \Gamma\vdash P:A\parr B
  \qquad
  \Delta\vdash Q:A^\bot
  \qquad
  \operatorname{TC}_{\parr}(P,Q)
}{
  \Gamma,\Delta\vdash P\,Q:B
}
\quad(\parr E).
\]
The side condition is the trace-class condition for the corresponding
evaluation pipeline.

\subsection{Linear negation and dual transposition}

Linear negation is interpreted by polarity reversal:
\[
\llbracket A^\bot\rrbracket=\llbracket A\rrbracket^\bot.
\]
A process
\[
F:\mathbb A\to\mathbb B
\]
determines a dual process
\[
F^\bot:\mathbb B^\bot\to\mathbb A^\bot
\]
by reversing polarities.  If
\[
F:
\mathbb A^+\times\mathbb B^-
\rightsquigarrow
\mathbb B^+\times\mathbb A^-,
\]
then
\[
F^\bot:
\mathbb B^-\times\mathbb A^+
\rightsquigarrow
\mathbb A^-\times\mathbb B^+
\]
is obtained by pre- and post-composition with the symmetry isomorphisms:
\[
F^\bot
:=
\sigma_{\mathbb B^+,\mathbb A^-}
\circ
F
\circ
\sigma_{\mathbb B^-,\mathbb A^+}.
\]
Equivalently,
\[
\begin{tikzcd}[column sep=huge,row sep=large]
\mathbb B^-\times\mathbb A^+
  \arrow[r, "\sigma_{\mathbb B^-,\mathbb A^+}"]
&
\mathbb A^+\times\mathbb B^-
  \arrow[r, "F"]
&
\mathbb B^+\times\mathbb A^-
  \arrow[r, "\sigma_{\mathbb B^+,\mathbb A^-}"]
&
\mathbb A^-\times\mathbb B^+ .
\end{tikzcd}
\]

The proof system has the involutive duality equations
\[
(A^\bot)^\bot=A,
\qquad
(M^\bot)^\bot=M,
\]
and
\[
(M\otimes N)^\bot=M^\bot\parr N^\bot,
\qquad
(M\parr N)^\bot=M^\bot\otimes N^\bot,
\]
interpreted by the canonical duality isomorphisms of the partial Int
paracategory.

\subsection{Cut, execution, and substitution}

The proof system has a general partial cut rule:
\[
\frac{
  \Gamma\vdash M:A
  \qquad
  \Delta,x:A\vdash N:B
  \qquad
  \operatorname{TC}_x(N,M)
}{
  \Delta,\Gamma\vdash N[M/x]:B
}
\quad(\operatorname{cut}).
\]
The side condition \(\operatorname{TC}_x(N,M)\) is the trace-class condition
for the open Geometry-of-Interaction pipeline obtained by connecting the
\(A\)-output of \(M\) to the \(A\)-input of \(N\).

Thus substitution is not a purely formal operation.  It is execution along a
feedback channel.  If the feedback equation has a unique internal solution,
the substitution is defined and is computed by substituting that solution into
the transition and output maps of the open pipeline.

\begin{theorem}[Soundness of multiplicative partial natural deduction]
\label{thm:pml-soundness}
Every derivation
\[
\Gamma\vdash M:A
\]
in the natural-deduction system above denotes a bidirectional Mealy process
\[
\llbracket M\rrbracket:
\llbracket\Gamma\rrbracket
\to
\llbracket A\rrbracket
\]
in
\[
\operatorname{Int}_{\partial}(\mathsf{Mly}(\mathcal E)).
\]
Every instance of cut, tensor elimination, unit elimination, implication
elimination, and par elimination is defined exactly when its associated
feedback equation is trace-class.
\end{theorem}

\begin{proof}
The proof is by induction on derivations.

The identity rule is interpreted by identity bidirectional Mealy processes.
Exchange is interpreted by symmetry isomorphisms.  Tensor introduction is
interpreted by the monoidal product of processes.  Unit introduction is
interpreted by the identity process on \(\mathbb I\).

The implication introduction rule is interpreted by the structural
Int-transpose of the denotation of the premise.  This transpose is total
because it is induced only by associativity, symmetry, unit, and polarity
rebracketing.  The par rule is the same construction under the definition
\(A\parr B=A^\bot\multimap B\).  Linear negation is interpreted by the
polarity-reversal operation on bidirectional interfaces and processes.

The elimination and substitution rules are interpreted by partial Int
composition.  Their side condition is precisely the condition that the open
pipeline being composed is trace-class.  When the side condition holds, the
partial composite exists by Theorem~\ref{thm:partial-goi-mealy}.  Hence every
derivable judgement has the displayed denotation.
\end{proof}

\begin{theorem}[Cut execution]
\label{thm:pml-cut-execution}
Let
\[
\Gamma\vdash M:A
\]
and
\[
\Delta,x:A\vdash N:B
\]
be derivations such that
\[
\operatorname{TC}_x(N,M)
\]
holds.  Then the denotation of
\[
N[M/x]
\]
is obtained by Geometry-of-Interaction execution: form the open pipeline,
solve the internal feedback equation uniquely over the cut interface
\[
\llbracket A\rrbracket,
\]
and substitute that solution into the transition and output maps of the
pipeline.
\end{theorem}

\begin{proof}
The denotation of \(N[M/x]\) is, by definition, the partial Int composite of
\[
\llbracket M\rrbracket
\]
and
\[
\llbracket N\rrbracket.
\]
By Definition~\ref{def:partial-int-composition-mealy}, this composite is the
partial trace of the corresponding open pipeline.  By
Proposition~\ref{prop:component-formula-partial-trace}, the transition and
output maps of this partial trace are obtained by substituting the unique
feedback solution into the transition and output maps of the open pipeline.
\end{proof}

\subsection{Orthogonality and behaviours}

The process calculus above has a standard GoI orthogonality layer.  This layer
turns collections of processes into types, following the use of orthogonality
in typed Geometry of Interaction
\cite{HaghverdiTypedGoI}.

Fix a class
\[
\mathcal S
\subseteq
\operatorname{Int}_{\partial}(\mathsf{Mly}(\mathcal E))(\mathbb I,\mathbb I)
\]
of successful closed processes.  Assume that \(\mathcal S\) is stable under
polarity identifications, tensorial reassociations, symmetries, and
closed-loop isomorphisms generated by defined partial Int composition.

For a bidirectional interface \(\mathbb A\), define the points and copoints
\[
\operatorname{Pt}(\mathbb A)
:=
\operatorname{Int}_{\partial}(\mathsf{Mly}(\mathcal E))(\mathbb I,\mathbb A),
\]
\[
\operatorname{CoPt}(\mathbb A)
:=
\operatorname{Int}_{\partial}(\mathsf{Mly}(\mathcal E))(\mathbb A,\mathbb I).
\]
There is a canonical polarity identification
\[
\operatorname{CoPt}(\mathbb A)
\cong
\operatorname{Pt}(\mathbb A^\bot),
\]
because an arrow
\[
\mathbb A\to\mathbb I
\]
is represented by a Mealy morphism
\[
\mathbb A^+\times\mathbf 1
\rightsquigarrow
\mathbf 1\times\mathbb A^-,
\]
which, after deleting terminal factors and reversing polarity, is precisely an
arrow
\[
\mathbb I\to\mathbb A^\bot.
\]

\begin{definition}[Trace-class orthogonality]
\label{def:trace-class-orthogonality}
For
\[
p\in\operatorname{Pt}(\mathbb A),
\qquad
q\in\operatorname{CoPt}(\mathbb A),
\]
write
\[
p\perp q
\]
if the partial composite
\[
q\circ_{\partial}p:\mathbb I\to\mathbb I
\]
is defined and belongs to \(\mathcal S\).
\end{definition}

For a set \(X\subseteq\operatorname{Pt}(\mathbb A)\), define
\[
X^\perp
:=
\{q\in\operatorname{CoPt}(\mathbb A)\mid
\forall p\in X,\ p\perp q\}.
\]
Using
\[
\operatorname{CoPt}(\mathbb A)
\cong
\operatorname{Pt}(\mathbb A^\bot),
\]
we regard \(X^\perp\) as a set of points of \(\mathbb A^\bot\).

For a set \(Y\subseteq\operatorname{CoPt}(\mathbb A)\), define
\[
Y^\perp
:=
\{p\in\operatorname{Pt}(\mathbb A)\mid
\forall q\in Y,\ p\perp q\}.
\]

\begin{definition}[Mealy conduct]
\label{def:mealy-conduct}
A \emph{Mealy conduct} on \(\mathbb A\) is a subset
\[
X\subseteq\operatorname{Pt}(\mathbb A)
\]
such that
\[
X=X^{\perp\perp}.
\]
\end{definition}

Conducts are the semantic propositions of the orthogonality model.  A process
\[
F:\mathbb A\to\mathbb B
\]
realizes a map of conducts
\[
X\to Y
\]
when, for every
\[
p\in X,
\]
the composite
\[
F\circ_{\partial}p
\]
is defined and belongs to \(Y\).

If \(X\) is a conduct on \(\mathbb A\) and \(Y\) is a conduct on
\(\mathbb B\), define
\[
X\otimes Y
:=
\{p\otimes q\mid p\in X,\ q\in Y\}^{\perp\perp},
\]
a conduct on
\[
\mathbb A\otimes\mathbb B.
\]
Linear negation is defined by orthogonality:
\[
X^\bot:=X^\perp,
\]
regarded as a conduct on
\[
\mathbb A^\bot.
\]
Par is defined by De Morgan duality:
\[
X\parr Y:=(X^\bot\otimes Y^\bot)^\bot,
\]
a conduct on
\[
\mathbb A\parr\mathbb B.
\]

\begin{theorem}[Conduct semantics for multiplicative partial linear logic]
\label{thm:conduct-semantics-partial-linear-logic}
With the stable success class \(\mathcal S\), Mealy conducts form the
multiplicative orthogonality model associated to
\[
\operatorname{Int}_{\partial}(\mathsf{Mly}(\mathcal E)).
\]
Tensor is interpreted by biorthogonal closure of parallel composition, linear
negation by trace-class orthogonality, par by De Morgan duality, and cut by
partial Int composition whenever the relevant feedback equation is
trace-class.
\end{theorem}

\begin{proof}
Biorthogonal closure gives stable semantic propositions.  Tensor of proofs is
parallel composition of bidirectional Mealy processes, followed by
biorthogonal closure at the level of conducts.  Negation is orthogonality,
transported along the canonical identification
\[
\operatorname{CoPt}(\mathbb A)
\cong
\operatorname{Pt}(\mathbb A^\bot).
\]
The cut rule is interpreted by partial composition in
\[
\operatorname{Int}_{\partial}(\mathsf{Mly}(\mathcal E)).
\]
When the cut is defined, the partial trace axioms give the usual invariance of
GoI execution under reassociation of cuts.  Orthogonality is preserved by the
assumed stability of \(\mathcal S\) under polarity identifications,
tensorial reassociations, symmetries, and closed-loop isomorphisms generated
by defined partial Int composition.  Therefore the multiplicative
orthogonality clauses are sound.
\end{proof}

\section{Execution formula for bidirectional processes}
\label{sec:execution-formula-bidirectional-processes}

The preceding construction admits a concrete execution formula.  Let
\[
F:\mathbb A\to\mathbb B,
\qquad
G:\mathbb B\to\mathbb C
\]
be bidirectional Mealy processes.  The open pipeline
\[
H_{G,F}
\]
has an internal feedback equation over the hidden interface \(\mathbb B^-\).
The partial composite exists exactly when this equation has a unique internal
solution.

\begin{definition}[Execution witness]
\label{def:execution-witness}
The \emph{execution witness object}
\[
W_{G,F}
\]
is the feedback witness object of
\[
H_{G,F}
\]
over \(\mathbb B^-\).  Its base object
\[
D_{G,F}
\]
is the object of visible admissible inputs and closed states for the open
pipeline \(H_{G,F}\).
\end{definition}

Thus \(W_{G,F}\) is defined by the instance
\[
W_{G,F}
=
W_{H_{G,F}},
\qquad
D_{G,F}
=
D_{H_{G,F}},
\]
of Definitions~\ref{def:feedback-witness-object} and
\ref{def:closed-state-object} applied to the Mealy morphism
\[
H_{G,F}:
\mathbb A^+\times\mathbb C^-\times\mathbb B^-
\rightsquigarrow
\mathbb C^+\times\mathbb A^-\times\mathbb B^-.
\]

\begin{proposition}[Execution criterion]
\label{prop:execution-criterion}
The partial composite
\[
G\circ_{\partial}F
\]
is defined if and only if
\[
W_{G,F}\to D_{G,F}
\]
is an isomorphism.
\end{proposition}

\begin{proof}
By Definition~\ref{def:partial-int-composition-mealy}, the partial composite
is defined exactly when the open pipeline \(H_{G,F}\) is trace-class over
\(\mathbb B^-\).  By
Proposition~\ref{prop:internal-trace-class-criterion}, this is equivalent to
the feedback witness projection
\[
W_{G,F}\to D_{G,F}
\]
being an isomorphism.
\end{proof}

\begin{theorem}[Mealy execution formula]
\label{thm:mealy-execution-formula}
Suppose
\[
G\circ_{\partial}F
\]
is defined.  Then its carrier is the closed-loop carrier of the open pipeline
\(H_{G,F}\), and its transition and output maps are obtained by substituting
the unique execution witness into the transition and output maps of
\(H_{G,F}\).
\end{theorem}

\begin{proof}
This is Proposition~\ref{prop:component-formula-partial-trace} applied to the
open pipeline \(H_{G,F}\).  The unique feedback solution is the execution
witness supplied by Proposition~\ref{prop:execution-criterion}.  Substituting
that witness into the transition and output maps of \(H_{G,F}\) gives the
traced Mealy morphism
\[
\operatorname{Tr}^{\mathbb B^-}(H_{G,F})
=
G\circ_{\partial}F.
\]
\end{proof}

\begin{remark}[Relation to Geometry of Interaction]
\label{rem:relation-goi}
The execution formula above is the Mealy analogue of the
Geometry-of-Interaction execution formula.  The categorical role of trace in
such execution formulas is classical in the traced-monoidal treatment of
Geometry of Interaction, for example in the work of
Abramsky--Haghverdi--Scott and Haghverdi
\cite{AbramskyHaghverdiScottGoI,HaghverdiTypedGoI}.  In the present setting
the execution formula is partial because the feedback witness must exist
uniquely in order for the result to remain Mealy.
\end{remark}

\section{Dynamic logic of bidirectional execution}
\label{sec:dynamic-logic-bidirectional-execution}

Let
\[
F:\mathbb A\to\mathbb B,
\qquad
G:\mathbb B\to\mathbb C
\]
be bidirectional processes whose partial composite is defined.  Let
\[
H_{G,F}
\]
be their open pipeline and let
\[
H_{G,F}^{\circ}
\]
be its closed-loop program category over the hidden interface \(\mathbb B^-\).

\begin{theorem}[Dynamic logic of Int-composition]
\label{thm:dynamic-logic-int-composition}
For every predicate \(\varphi\) on the state object of
\[
G\circ_{\partial}F,
\]
one has
\[
\langle
\mathbf m_{G\circ_{\partial}F}
\rangle\varphi
=
\langle
\mathbf m_{H_{G,F}}^{\circ}
\rangle\varphi.
\]
In the first-order Heyting setting,
\[
[
\mathbf m_{G\circ_{\partial}F}
]\varphi
=
[
\mathbf m_{H_{G,F}}^{\circ}
]\varphi.
\]
\end{theorem}

\begin{proof}
The partial composite is the partial trace of \(H_{G,F}\).  Therefore its
program category is canonically the closed-loop program category of
\(H_{G,F}\), by Theorem~\ref{thm:program-category-of-traced-mealy}.  The
diamond and box predicate transformers are invariant under the resulting
isomorphism of primitive execution spans.
\end{proof}

\section{Representable systems and trace}
\label{sec:representable-systems-and-trace}

We record the interaction between trace and the stateless representable
fragment.

Let
\[
F:\mathbb A\times\mathbb U\to\mathbb B\times\mathbb U
\]
be an internal functor, regarded as a representable Mealy morphism.  Its
representable carrier is
\[
P=A_0\times U_0,
\]
with left leg the identity on \(A_0\times U_0\), and right leg the object map
of \(F\):
\[
\begin{tikzcd}[column sep=large]
A_0\times U_0
&
A_0\times U_0
  \arrow[l, "1"']
  \arrow[r, "F_0"]
&
B_0\times U_0 .
\end{tikzcd}
\]

\begin{proposition}[Trace of a representable Mealy morphism]
\label{prop:trace-representable-mealy}
If the representable Mealy morphism induced by \(F\) is trace-class over
\(\mathbb U\), then
\[
\operatorname{Tr}^{\mathbb U}(F)
\]
is a Mealy morphism
\[
\mathbb A\rightsquigarrow\mathbb B.
\]
It is representable precisely when the closed carrier object of the trace is
isomorphic over \(A_0\) to \(A_0\).
\end{proposition}

\begin{proof}
Trace-class gives a traced Mealy morphism by
Theorem~\ref{thm:partial-trace-on-mealy}.  The carrier of the trace is the
closed state object
\[
P^\circ.
\]
A Mealy morphism is representable precisely when its carrier span has
invertible left leg.  In the present case this left leg is
\[
P^\circ\to A_0.
\]
Hence the trace is representable precisely when this map is an isomorphism.
\end{proof}

\begin{remark}[Tracing can create state]
\label{rem:tracing-can-create-state}
Even when the open system is induced by an internal functor, feedback closure
may produce a genuinely stateful Mealy morphism.  It remains an internal
functor only when the feedback closure introduces no hidden state, equivalently
when the closed carrier is isomorphic to \(A_0\) over \(A_0\).
\end{remark}

\section{The one-object case}
\label{sec:one-object-partial-trace}

Let
\[
\mathcal E=\mathbf{Set},
\]
and suppose that
\[
\mathbb A=M,
\qquad
\mathbb B=N,
\qquad
\mathbb U=K
\]
are one-object internal categories, hence monoids.  A Mealy morphism
\[
M\times K\rightsquigarrow N\times K
\]
has state set \(P\) and functions
\[
\delta:(M\times K)\times P\to P,
\]
\[
\omega_N:(M\times K)\times P\to N,
\qquad
\omega_K:(M\times K)\times P\to K.
\]

\begin{proposition}[One-object trace criterion]
\label{prop:one-object-trace-criterion}
The Mealy morphism above is trace-class over \(K\) precisely when, for every
\[
m\in M
\]
and every state \(x\in P\), the equation
\[
k=\omega_K(m,k,x)
\]
has a unique solution \(k\in K\).
When this holds, the traced Mealy morphism
\[
M\rightsquigarrow N
\]
is given by
\[
\delta^{\operatorname{tr}}(m,x)
=
\delta(m,k(m,x),x),
\]
and
\[
\omega^{\operatorname{tr}}(m,x)
=
\omega_N(m,k(m,x),x),
\]
where \(k(m,x)\) is the unique solution.
\end{proposition}

\begin{proof}
This is the specialization of
Proposition~\ref{prop:internal-trace-class-criterion} and
Proposition~\ref{prop:component-formula-partial-trace} to one-object internal
categories.  The closed-state condition is automatic in the one-object case.
\end{proof}

\begin{remark}[Convention on multiplication order]
\label{rem:one-object-trace-order}
The displayed one-object formulas follow the multiplication and composition
order fixed in the earlier chapters.  Readers using the opposite convention
for monoid multiplication should reverse the relevant composite orders.
\end{remark}

\begin{remark}[Genuine partiality]
\label{rem:genuine-partiality-one-object}
If
\[
\omega_K(m,k,x)=k
\]
for all \(k\), then the feedback equation has many solutions unless \(K\) is
terminal.  If \(K\) is a group and
\[
\omega_K(m,k,x)=kc
\]
for a nontrivial \(c\in K\), then the equation
\[
k=kc
\]
has no solution.  Thus the trace is not total in general.
\end{remark}

\section{Main consolidation}
\label{sec:partial-trace-main-consolidation}

\begin{theorem}[Partial trace, feedback, privacy, and bidirectional Mealy processes]
\label{thm:partial-trace-main-consolidation}
The constructions of this chapter give the following hierarchy.

\begin{enumerate}
\item Mealy morphisms form a symmetric monoidal category
\[
\mathsf{Mly}(\mathcal E)
\]
under product of internal categories.

\item The labelled input-discrete presentation embeds
\[
\mathsf{Mly}(\mathcal E)
\]
as the symmetric monoidal subcategory of
\[
\mathbf{Span}(\mathbf{Cat}_{\mathrm{int}}(\mathcal E))
\]
whose morphisms are spans of internal categories with input-discrete left leg.

\item The ambient span category carries the standard total trace.

\item A Mealy morphism
\[
\Phi:\mathbb A\times\mathbb U
\rightsquigarrow
\mathbb B\times\mathbb U
\]
is trace-class over \(\mathbb U\) exactly when the ambient traced span remains
input-discrete.

\item Equivalently, if
\[
\Phi=(P,\delta,\omega),
\]
then \(\Phi\) is trace-class precisely when the internal feedback equation
\[
h=\omega_U((a,h),x)
\]
has a unique solution for every closed state \(x\in P^\circ\) and every
visible input arrow \(a\) admissible at \(x\).

\item The trace
\[
\operatorname{Tr}^{\mathbb U}(\Phi):
\mathbb A\rightsquigarrow\mathbb B
\]
has carrier
\[
\begin{tikzcd}[column sep=large]
A_0
&
P^\circ
  \arrow[l, "p_A"']
  \arrow[r, "q_B"]
&
B_0
\end{tikzcd}
\]
and transition-output maps obtained by substituting the unique feedback
solution.

\item Therefore
\[
\mathsf{Mly}(\mathcal E)
\]
is a partially traced symmetric monoidal category in the sense of
Malherbe--Scott--Selinger.

\item Tracing over \(\mathbb U\) has an interface-privacy interpretation: the
feedback witness is used internally but hidden from the visible
\(\mathbb A\)-to-\(\mathbb B\) labelled behaviour.

\item The program category of the traced Mealy morphism is the closed-loop
program category:
\[
\mathsf{Prog}_{\operatorname{Tr}^{\mathbb U}(\Phi)}
\cong
\operatorname{Eq}(I_{\Phi,U},\Omega_{\Phi,U}).
\]

\item Consequently, the diamonds and boxes of the traced system are exactly
the diamonds and boxes of the closed-loop primitive execution span.

\item Relative semantics commutes with feedback closure:
\[
\mathsf{Prog}_{\operatorname{Tr}^{\mathbb U}(\Phi),Y}
\cong
\left(
\mathsf{Prog}_{\Phi,Y\times\mathbf 1_{\mathbb U}}
\right)^\circ.
\]

\item Response coalgebras carry a partial feedback trace.  For Mealy-induced
coalgebras,
\[
\operatorname{Tr}^{\mathbb U}(c_\Phi)
=
c_{\operatorname{Tr}^{\mathbb U}(\Phi)}.
\]

\item Flattening commutes with feedback trace:
\[
\operatorname{Flat}
\left(
\operatorname{Tr}^{\mathbb U}(c_\Phi)
\right)
\cong
\mathbf m_{\Phi}^{\circ}.
\]

\item The partial Int-construction yields a symmetric monoidal paracategory
\[
\operatorname{Int}_{\partial}(\mathsf{Mly}(\mathcal E))
\]
of bidirectional Mealy processes.

\item Partial composition in this paracategory is Geometry-of-Interaction
execution: form the open pipeline and partially trace the shared counterflow
interface.

\item The internal proof system of this partial Geometry-of-Interaction
paracategory is multiplicative partial linear logic.  Formulas are
bidirectional Mealy interfaces, proofs are bidirectional Mealy processes, and
cut is defined exactly when the corresponding open pipeline is trace-class.

\item With trace-class orthogonality and a stable class
\(\mathcal S\) of successful closed executions, Mealy conducts give the
corresponding GoI orthogonality model: tensor is parallel composition, linear
negation is orthogonality, and cut elimination is unique feedback substitution.
\end{enumerate}
\end{theorem}

\begin{proof}
Part (1) is Proposition~\ref{prop:mly-symmetric-monoidal}.  Part (2) is
Theorem~\ref{thm:mealy-as-span-subcategory}.  Part (3) is
Proposition~\ref{prop:ambient-span-total-trace}.  Part (4) is
Proposition~\ref{prop:trace-class-ambient-closure}.  Part (5) is
Proposition~\ref{prop:internal-trace-class-criterion}.  Part (6) is
Proposition~\ref{prop:component-formula-partial-trace}.  Part (7) is
Theorem~\ref{thm:partial-trace-on-mealy}.  Part (8) is
Proposition~\ref{prop:privacy-by-feedback-hiding}.  Part (9) is
Theorem~\ref{thm:program-category-of-traced-mealy}.  Part (10) is
Corollary~\ref{cor:diamonds-boxes-feedback-closure}.  Part (11) is
Theorem~\ref{thm:relative-semantics-commutes-feedback}.  Part (12) is
Proposition~\ref{prop:trace-mealy-induced-coalgebras}.  Part (13) is
Theorem~\ref{thm:flattening-commutes-feedback}.  Parts (14) and (15) are
Theorem~\ref{thm:partial-goi-mealy} and
Theorem~\ref{thm:mealy-execution-formula}.  Parts (16) and (17) are
Theorem~\ref{thm:pml-soundness}, Theorem~\ref{thm:pml-cut-execution}, and
Theorem~\ref{thm:conduct-semantics-partial-linear-logic}.
\end{proof}

\begin{remark}[Conceptual summary]
\label{rem:partial-trace-conceptual-summary}
The partial trace structure is not imposed on Mealy morphisms.  It emerges
because Mealy morphisms are precisely those labelled spans of internal
categories whose left leg is input-discrete.  The total span trace closes a
hidden interface by equalizing hidden input and hidden output.  This closure
remains Mealy exactly when input-discreteness survives.  Internally, that
condition is the unique-solvability of the feedback equation
\[
h=\omega_U((a,h),x).
\]
Thus partial trace is deterministic feedback in the Mealy calculus.  Its
bidirectional Int-construction gives Geometry-of-Interaction execution, and
its internal proof theory is the multiplicative partial linear logic whose cut
rule is governed by trace-class feedback.
\end{remark}